\documentclass[11pt]{article}

\usepackage[margin=1.1in]{geometry}

\usepackage{amssymb,mathtools,natbib}

% Anchored formal environments for causalsmith papers.
% First mandatory argument = obj_id (machine anchor joining the paper to the
% Lean crosswalk; invisible in the rendered PDF). Optional argument = name.
% Each environment also defines \label{obj:<obj_id>} so prose can use
% \cref{obj:T-1}; cleveref derives the printed kind and number from the target.
\usepackage{amsmath}
\usepackage{amsthm}
% Long displays may break across pages; let TeX stretch a little harder before an
% overfull \hbox so wide formulas/tables are less likely to run off the page.
\allowdisplaybreaks
\setlength{\emergencystretch}{3em}
\newtheorem{theorem}{Theorem}
\newtheorem{lemma}{Lemma}
\newtheorem{assumption}{Assumption}
\newtheorem{definition}{Definition}
\newtheorem{proposition}{Proposition}
\newtheorem{remark}{Remark}
\newtheorem{citedresult}{Cited result}
\NewDocumentEnvironment{theoremv}{m o s}{\begin{theorem}\label{obj:#1}{\normalfont[\texttt{\detokenize{#1}}]\IfValueT{#2}{\ (#2)}\IfBooleanT{#3}{\verificationfootnotemark}\ }}{\end{theorem}}
\NewDocumentEnvironment{lemmav}{m o s}{\begin{lemma}\label{obj:#1}{\normalfont[\texttt{\detokenize{#1}}]\IfValueT{#2}{\ (#2)}\IfBooleanT{#3}{\verificationfootnotemark}\ }}{\end{lemma}}
\NewDocumentEnvironment{assumptionv}{m o s}{\begin{assumption}\label{obj:#1}{\normalfont[\texttt{\detokenize{#1}}]\IfValueT{#2}{\ (#2)}\IfBooleanT{#3}{\verificationfootnotemark}\ }}{\end{assumption}}
\NewDocumentEnvironment{definitionv}{m o s}{\begin{definition}\label{obj:#1}{\normalfont[\texttt{\detokenize{#1}}]\IfValueT{#2}{\ (#2)}\IfBooleanT{#3}{\verificationfootnotemark}\ }}{\end{definition}}
% A `propositionv` is a proved result tiered as a Proposition (theorem-grade, machine-verified).
\NewDocumentEnvironment{propositionv}{m o s}{\begin{proposition}\label{obj:#1}{\normalfont[\texttt{\detokenize{#1}}]\IfValueT{#2}{\ (#2)}\IfBooleanT{#3}{\verificationfootnotemark}\ }}{\end{proposition}}
% A `remarkv` holds NON-load-bearing interpretive / open-question content (e.g. a causalsmith `oeq:`
% open question). It is neither proved nor assumed; no theorem may depend on it.
\NewDocumentEnvironment{remarkv}{m o s}{\begin{remark}\label{obj:#1}{\normalfont[\texttt{\detokenize{#1}}]\IfValueT{#2}{\ (#2)}\IfBooleanT{#3}{\verificationfootnotemark}\ }}{\end{remark}}
% A `citedv` is an IMPORTED external result the paper relies on but does NOT prove
% (a causalsmith `gate_class:"cited"` node). It is machine-checked against the cited
% source at F2.5, never proved here; the fixed provenance note makes that explicit
% so the environment can never be read as a contribution of this paper. The \citep
% attribution is supplied by the surrounding prose (P2), keyed to references.bib.
\NewDocumentEnvironment{citedv}{m o s}{\begin{citedresult}\label{obj:#1}{\normalfont[\texttt{\detokenize{#1}}]\IfValueT{#2}{\ (#2)}\IfBooleanT{#3}{\verificationfootnotemark}\ \textit{(imported result; machine-checked against the cited source, not proved here.)}\ }}{\end{citedresult}}
% \leanref{obj_id}{display text}: an inline first-mention link to a from-note object's
% verified Lean code. In the PDF it renders the display text plainly (the Lean link is a
% web-only affordance); tex2html turns it into a drawer-opening span.
\newcommand{\leanref}[2]{#2}
% For a theorem/lemma whose Lean declaration accepts a source-matched published proposition as an
% input, the starred anchored environment puts the mark in its heading; the generated text remains
% outside the frozen mathematical environment. Keep the combined macro for older paper bundles.
\newcommand{\verificationfootnotemark}{\footnotemark}
\newcommand{\verificationfootnotetext}[1]{\footnotetext{#1}}
\newcommand{\verificationfootnote}[1]{\verificationfootnotemark\verificationfootnotetext{#1}}
% hyperref follows natbib and the environment macros; cleveref must follow hyperref.
\usepackage[colorlinks=true,linkcolor=blue,citecolor=blue,urlcolor=blue]{hyperref}
\usepackage[nameinlink,noabbrev]{cleveref}
% Pin the names explicitly: labels are placed inside the underlying amsthm environments, and these
% declarations make the PDF contract independent of cleveref's theorem-name inference.
\crefname{theorem}{theorem}{theorems}
\Crefname{theorem}{Theorem}{Theorems}
\crefname{lemma}{lemma}{lemmas}
\Crefname{lemma}{Lemma}{Lemmas}
\crefname{assumption}{assumption}{assumptions}
\Crefname{assumption}{Assumption}{Assumptions}
\crefname{definition}{definition}{definitions}
\Crefname{definition}{Definition}{Definitions}
\crefname{proposition}{proposition}{propositions}
\Crefname{proposition}{Proposition}{Propositions}
\crefname{remark}{remark}{remarks}
\Crefname{remark}{Remark}{Remarks}
\crefname{citedresult}{cited result}{cited results}
\Crefname{citedresult}{Cited result}{Cited results}


% Permit inline math delimiters in frozen formal statements that place prose in
% display math; ordinary inline uses retain their usual behavior.
\newif\ifcausalsmithinlineopened
\renewcommand{\(}{\ifmmode\causalsmithinlineopenedfalse\else\causalsmithinlineopenedtrue\begin{math}\fi}
\renewcommand{\)}{\ifcausalsmithinlineopened\end{math}\causalsmithinlineopenedfalse\fi}

\title{Population Identification of Competing Latent LiNGAM Representations by Higher-Order Cumulants}

\date{\today}

\begin{document}

\maketitle

\begin{abstract}
We study a population identification question for a bivariate latent linear non-Gaussian model in which the number \(m\) of middle source slots and an axis normalization are maintained inputs. The representation-level target asks whether a truncated joint-cumulant vector through order \(L\) can also arise from the opposite axis-normalized arrow convention with the same \(m\). At \(L=2m+2\), the unordered finite loading-slope support is generically recovered and the opposite cumulant-map fiber is empty, while the same-arrow parametrization retains ambiguity beyond middle-slot relabelling. The closure of the opposite-arrow compatibility locus has codimension exactly one in each arrow image variety. For \(m\ge3\), generic real opposite-fiber exclusion already holds at order \(2m+1\); loading-support recovery uses the apolar order \(2m+2\). The real-incidence result applies to Euclidean-open parameter neighborhoods meeting the feasible region and sharpens the population geometry of the maintained representation class.
\end{abstract}

\section{Introduction}

Econometric models often become directionally informative only after the researcher imposes structure beyond covariance restrictions. In linear non-Gaussian models, higher-order cumulants encode the loading directions of independent shocks and can therefore restrict which recursive representation is compatible with the population distribution. This mechanism underlies LiNGAM and its latent-variable extensions \citep{ShimizuHyvarinenKanoHoyer2006,HoyerShimizuKerminenPalviainen2008,TashiroShimizuHyvarinenWashio2014,MaedaShimizu2020}, and is closely connected to independent component analysis and cumulant-based tensor decompositions \citep{Comon1994,HyvarinenOja2000,Cardoso1999,KoldaBader2009}.

This paper isolates a population question for a fixed number \(m\) of known or externally specified middle source slots and the two axis-normalized bivariate latent linear non-Gaussian representation classes in \cref{obj:def:forward-lvlingam-class} and~\cref{obj:def:reverse-lvlingam-class}. One convention tags the representation as \(X\to Y\), the other as \(Y\to X\). The observed object is the finite vector of joint cumulants \(T_L(P)\) in \cref{obj:def:truncated-cumulant}. The target is whether this vector has an opposite-arrow cumulant-map preimage with the same source count and normalization. It serves settings in which an institutional model, factor design, or separate model-selection analysis fixes the number of common shock slots and the researcher evaluates population compatibility of the maintained reverse representation.

\paragraph{Econometric target and decision.}
The maintained model class contains two competing structural representations with the same externally fixed \(m\). The population decision rejects a convention when the observed truncated cumulant vector has an empty fiber under its polynomial cumulant map. The resulting exclusion restriction is representation-level and comes with an exact description of the compatibility set and the residual ambiguity within the surviving convention. These questions matter before estimation: they establish the population target for a cumulant-based procedure and identify the algebraic surface near which such a procedure would be weak.

\paragraph{Contributions.}
The paper establishes three facts about this target. First, the apolar order \(2m+2\) generically recovers the loading-slope support and excludes the opposite fiber. Second, opposite-arrow compatibility has a codimension-one closure, substantially larger than a dimension count based on the \(m\) middle slots might suggest. Third, for \(m\ge3\), generic real opposite-fiber exclusion requires one less cumulant order than support recovery. A separate fiber result characterizes the remaining same-arrow parameter ambiguity.

The main result at the apolar order \(K=2m+2\) is \cref{obj:thm:generic-apolar-arrow-recovery}. On relatively Zariski-open dense parameter loci, and on Euclidean-open parameter neighborhoods that meet the real feasible region, the order-\(K\) cumulants recover the unordered finite loading-slope support through binary-form apolar equations and rule out a full representation in the opposite arrow convention. Its exact domain is the pair of polynomial cumulant maps and their parameter fibers.

The same-arrow corollary in \cref{obj:thm:generic-arrow-recovery-and-fiber-obstruction} adds the corresponding fiber behavior on the generic loci of the apolar theorem. The same cumulants preserve the unordered loading-slope multiset while leaving multiple same-arrow parametrizations beyond the admissible middle-source relabellings in \cref{obj:def:admissible-source-swaps}. In particular, a direct slot can be exchanged with a middle slot while preserving the cumulant vector and remaining outside the admissible label orbit. Thus generic arrow separation delivers more information than covariance equivalence while retaining precisely characterized structural-parameter ambiguity.

The third result, \cref{obj:thm:exceptional-locus-codimension-one}, describes the failure set at the same apolar order. The cumulant vectors that have a generic representation for one arrow and at least one full representation for the other form an exceptional compatibility locus whose closure has codimension exactly one inside each arrow image variety. For \(m\ge2\), a single algebraic dimension is therefore lost throughout this compatibility geometry. This codimension-one statement is important for interpretation: generic exclusion holds off the exceptional set, whose size makes proximity to compatibility surfaces empirically relevant.

The final identification result concerns the cumulant information needed for generic real opposite-arrow exclusion. \Cref{obj:thm:improved-real-information-order} shows that, for \(m\ge3\), order \(2m+1\) already suffices for generic separation over the real feasible regions of \cref{obj:def:real-feasible-regions}. This improves the real information-order bound relative to the apolar support-recovery order, which remains the relevant threshold when the target includes recovery of the loading support.

The paper is closest to the cumulant-based latent LiNGAM literature. \citet{ChenPengHuangEtAl2025Direction} develop a directional testing construction from a higher-cumulant rank deficiency at order \(2m+3\). \citet{SchkodaRobevaDrton2024} give recursive higher-cumulant rank conditions for unobserved confounding. \citet{TramontanoKivvaSalehkaleybarDrtonKiyavash2024,TramontanoKivvaSalehkaleybarDrtonKiyavash2025} target causal-effect identification from cumulants. Our one-order improvement pertains specifically to fixed-slot, axis-normalized population fiber exclusion; the comparison below aligns each result with its statistical or algebraic target.

The comparison with \citet{ChenPengHuangEtAl2025Direction} turns on distinct targets. In their result, \(m\) counts latent confounders in a real statistical testing problem, and the order-\(2m+3\) rank condition supports a causal-direction test. Here \(m\) counts externally fixed middle source slots, the axis normalization is maintained, and the conclusion is generic emptiness of an opposite polynomial-map fiber together with support-recovery and codimension-one geometry. The difference in order is therefore indexed to these respective targets and assumptions.

\begin{table}[ht]
\centering
\scriptsize
\caption{Auditable comparison with the closest cumulant-based approaches}
\begin{tabular}{p{0.17\linewidth}p{0.15\linewidth}p{0.17\linewidth}p{0.21\linewidth}p{0.20\linewidth}}
\hline
Work & Source-count and normalization & Population object & Scope of conclusion & Statistical output \\
\hline
Classical LiNGAM/ICA & model-specific dimension and mixing structure & observational distribution or mixing representation & ordering or mixing identification under the stated statistical model & estimators or algorithms are central \\
\citet{ChenPengHuangEtAl2025Direction} & \(m\) enters a latent-confounder testing construction & higher-cumulant rank condition & directional testing at order \(2m+3\) in a real statistical model & test construction \\
\citet{SchkodaRobevaDrton2024} & latent-confounding structure encoded by recursive rank restrictions & higher-cumulant tensors & rank conditions for unobserved confounding & identification conditions \\
\citet{TramontanoKivvaSalehkaleybarDrtonKiyavash2024,TramontanoKivvaSalehkaleybarDrtonKiyavash2025} & model-specific latent structure & cumulants relevant to causal effects & causal-effect identification & constructive identification methodology \\
This paper & externally fixed \(m\), fixed source slots, and axis normalization & truncated population cumulant vector & complex-generic opposite-fiber exclusion and support recovery at \(2m+2\); codimension-one compatibility; separate real-generic exclusion at \(2m+1\) for \(m\ge3\) & population fiber characterization \\
\hline
\end{tabular}
\end{table}

\paragraph{A one-middle-slot interpretation.}
When \(m=1\), the forward convention can be written
\[
X=S_0+S_1,\qquad
Y=\gamma S_0+\rho S_1+S_2.
\]
An econometrician may regard \(S_1\) as one externally maintained common-shock slot and \(S_0,S_2\) as the two axis-normalized innovation slots. If both loadings of \(S_1\) are nonzero, it has the usual latent-confounder interpretation; the theorem uses the more general source-slot representation. At order four, \cref{obj:thm:generic-apolar-arrow-recovery} says that, off its stated algebraic exceptional set, the resulting population cumulants have an empty reverse-convention fiber under the same one-slot specification. \Cref{obj:thm:exceptional-locus-codimension-one} then places vectors admitting both conventions on a codimension-one closure. This example illustrates the precise representation-level decision delivered by the theory.

The paper's scope is population identification from finite cumulants within a fixed axis-normalized latent-source representation class. It establishes representation compatibility and separation for a maintained \(m\), complementing finite-sample inference, model selection, and broader graphical identification analyses \citep{Pearl2009,SpirtesGlymourScheines2001,PetersJanzingScholkopf2017,RichardsonSpirtes2002}. This precise scope makes the causal interpretation conditional on the stated latent-source class.

The proofs use classical apolarity of binary forms and simultaneous decomposition ideas \citep{ComonMourrain1996,Landsberg2012}. Cumulants turn independent-source representations into sums of powers of loading directions; apolar equations then recover, or partially constrain, the support of those directions. The verification appendix records the Lean 4 declarations corresponding to the delivered mathematical objects.

\section{Setup and assumptions}

\subsection{The algebraic theorem domain}

The main theorems concern two polynomial maps defined on fixed-source-count parameter spaces. Fix \(m\) and a truncation order \(L\). The right-arrow parameter records the finite slopes \(\gamma,\rho_1,\ldots,\rho_m\) and order-specific source weights \(c_{jr}\); the left-arrow parameter records \(\delta,\sigma_1,\ldots,\sigma_m\) and \(d_{jr}\). For each \(2\le r\le L\), the two maps send these coordinates to the mixed-cumulant block
\[
t_{r,a}
=\sum_{j=0}^{m+1}c_{jr}(u_{j1})^{r-a}(u_{j2})^a
\quad\text{or}\quad
t_{r,a}
=\sum_{j=0}^{m+1}d_{jr}(v_{j1})^{r-a}(v_{j2})^a,
\qquad 0\le a\le r.
\]
The loading vectors \(u_j\) and \(v_j\) encode the two axis conventions. The theorem domain is the resulting pair of complex affine parameter spaces and polynomial cumulant maps, stated formally in \cref{obj:def:truncated-cumulant}--\cref{obj:def:reverse-cumulant-map}. ``Generic'' in the complex results means outside a proper algebraic subset of these parameter spaces. The conclusion establishes emptiness of one map fiber within this domain.

\subsection{Probabilistic interpretation}

The probabilistic assumptions now attach an econometric interpretation to real points of the algebraic maps. Each representation has \(m+2\) source slots: two are tied to the coordinate-axis normalization and the remaining \(m\) middle source slots have finite slopes. Centered mutually independent sources make cumulants additive across slots, finite moments make the retained cumulants well-defined, and non-Gaussianity provides higher-order information. These assumptions define the real feasible regions below, while the complex-polynomial interpretation of the maps applies on the full affine parameter spaces. A middle source slot receives the latent-confounder interpretation when both coordinate loadings are nonzero, so ``middle source slot'' is the default term.

The source dimension is indexed by \leanref{sym:m}{$m$}, with \leanref{sym:n}{$n$} later denoting \(m+2\) loading directions. The working truncation order \leanref{sym:K}{$K$} is finite, and \leanref{sym:S_j}{$S_j$} denotes the \(j\)th centered latent source.

\begin{assumptionv}{ass:independent-sources}[Independent latent sources]
The latent source variables \(S_0,\ldots,S_{m+1}\) are mutually independent with respect to \(\mu\):
\[
S_0\perp\cdots\perp S_{m+1}.
\]
\end{assumptionv}

\cref{obj:ass:independent-sources} is the standard independent-component source condition in latent-variable LiNGAM; it is what makes joint cumulants additive across source slots \citep{ChenPengHuangEtAl2025Direction}.

\begin{assumptionv}{ass:finite-cumulants}
For every \(j=0,\ldots,m+1\), the latent source variable \(S_j\) has finite \(K\)-th absolute moment:
\[
\mathbb E |S_j|^K<\infty .
\]
\end{assumptionv}

\cref{obj:ass:finite-cumulants} is the standard finite-order cumulant existence condition, tailored to the truncation order used by the paper \citep{ChenPengHuangEtAl2025Direction}.

\begin{assumptionv}{ass:source-nongaussianity}
For every \(j=0,\ldots,m+1\), the latent source variable \(S_j\) is non-Gaussian.
\end{assumptionv}

\cref{obj:ass:source-nongaussianity} is the usual non-Gaussian disturbance condition in LiNGAM-type identification; higher-order non-Gaussian structure supplies the cumulant information beyond covariance restrictions \citep{ChenPengHuangEtAl2025Direction}.

The two axis-normalized parametrizations give the real structural interpretation of the arrow tags used by the polynomial maps. In the right-arrow convention, \leanref{sym:X}{$X$} is the first coordinate and \leanref{sym:Y}{$Y$} is the second coordinate, the direct loading has slope \leanref{sym:gamma}{$\gamma$}, the middle finite slopes are \leanref{sym:rho_i}{$\rho_i$}, and the forward loading vectors are \leanref{sym:u_j}{$u_j$}. In the left-arrow convention, the analogous direct and middle slopes are \leanref{sym:delta}{$\delta$} and \leanref{sym:sigma_i}{$\sigma_i$}, with reverse loading vectors \leanref{sym:v_j}{$v_j$}.

\begin{assumptionv}{ass:forward-axis-model}
With forward loading vectors \(u_0=(1,\gamma)\), \(u_j=(1,\rho_j)\) for \(1\le j\le m\), and \(u_{m+1}=(0,1)\), the bivariate observed outcome vector admits the forward latent linear representation
\[
(X,Y)^\top=\sum_{j=0}^{m+1} u_j S_j .
\]
\end{assumptionv}

\cref{obj:ass:forward-axis-model} is a bivariate linear acyclic representation with latent source slots in the \(X\to Y\) orientation \citep{ChenPengHuangEtAl2025Direction}. The axis normalization fixes one horizontal-axis and one vertical-axis slot, so only the middle source slots carry free finite slopes.

\begin{assumptionv}{ass:reverse-axis-model}
With reverse loading vectors \(v_0=(1,0)\), \(v_j=(\sigma_j,1)\) for \(1\le j\le m\), and \(v_{m+1}=(\delta,1)\), the bivariate observed outcome vector admits the reverse latent linear representation
\[
(X,Y)^\top=\sum_{j=0}^{m+1} v_j S_j .
\]
\end{assumptionv}

\cref{obj:ass:reverse-axis-model} is the corresponding reversed bivariate linear acyclic parametrization for the \(Y\to X\) orientation \citep{ChenPengHuangEtAl2025Direction}. Its symmetric normalization permits a direct comparison of the two representation fibers.

Genericity enters through distinct finite loading directions and nonzero direct edges. These restrictions remove lower-dimensional coincidences in the loading support while permitting middle source slots with a zero coordinate loading.

\begin{assumptionv}{ass:forward-noncollinearity}
The direct and latent finite slopes in the forward parametrization are pairwise distinct:
\[
\bigl|\{\gamma,\rho_i:1\le i\le m\}\bigr|=m+1 .
\]
\end{assumptionv}

\cref{obj:ass:forward-noncollinearity} is the standard generic distinct-loading-directions condition for the forward parametrization \citep{SchkodaRobevaDrton2024}. It excludes merged finite slopes while permitting zero-valued middle slopes.

\begin{assumptionv}{ass:reverse-noncollinearity}
The reverse slopes are pairwise distinct:
\[
\left|\{\delta,\sigma_i:1\le i\le m\}\right|=m+1.
\]
\end{assumptionv}

\cref{obj:ass:reverse-noncollinearity} imposes the same generic distinct-loading-directions condition on the reverse parametrization \citep{SchkodaRobevaDrton2024}. The symmetry is useful because later separation statements compare full opposite-arrow fibers.

\begin{assumptionv}{ass:forward-nonzero-edge}
The forward direct slope is nonzero:
\[
\gamma\ne 0.
\]
\end{assumptionv}

\cref{obj:ass:forward-nonzero-edge} is the standard nonzero direct-edge genericity restriction in the forward direction \citep{ChenPengHuangEtAl2025Direction}. It excludes the degenerate case in which the distinguished direct slot lies on the horizontal axis.

\begin{assumptionv}{ass:reverse-nonzero-edge}
The reverse direct slope is nonzero:
\[
\delta\ne 0.
\]
\end{assumptionv}

\cref{obj:ass:reverse-nonzero-edge} is the analogous nonzero direct-edge genericity restriction for the reverse direction \citep{ChenPengHuangEtAl2025Direction}. Together, the two nonzero-edge restrictions keep the two arrow conventions from sharing their distinguished direct-axis slot.

We now define the two law classes. The notation \(\operatorname{Laws}(\mathbb R^2)\) denotes the ambient class of Borel probability laws on \(\mathbb R^2\), and \leanref{S-1}{the observational law setup} is the data-generating world in which these law classes live.

\begin{definitionv}{def:forward-lvlingam-class}[Forward class \(\mathcal P^{\mathrm{right}}_{m,K}\)]
Write \(\operatorname{Laws}(\mathbb R^2)\) for the Borel probability laws on \(\mathbb R^2\). Let \(\mathcal P^{\mathrm{right}}_{m,2m+2}\) be the set of all probability laws \(P\in \operatorname{Laws}(\mathbb R^2)\) for which there exist a measurable space \((\Omega,\mathcal A)\), a probability measure \(\mu\) on \(\Omega\), source functions
\[
S_0,\ldots,S_{m+1}:\Omega\to\mathbb R,
\]
observed variables \(X,Y:\Omega\to\mathbb R\), a direct forward slope \(\gamma\in\mathbb R\), and latent forward slopes \(\rho_1,\ldots,\rho_m\in\mathbb R\) such that:
\begin{itemize}
\item \textbf{(Source independence.)} The sources \(S_0,\ldots,S_{m+1}\) are mutually independent under \(\mu\).
\item \textbf{(Finite cumulant order.)} Each source has finite moments through order \(2m+2\).
\item \textbf{(Source non-Gaussianity.)} Each source has non-Gaussian law under \(\mu\).
\item \textbf{(Centering.)} Each source is centered: \(\int S_j\,d\mu=0\) for every \(j=0,\ldots,m+1\).
\item \textbf{(Forward representation.)} For every \(\omega\in\Omega\), the bivariate observed outcome vector \((X(\omega),Y(\omega))^\top\) admits the forward \(X\to Y\) latent linear non-Gaussian representation with loading vectors \(u_j\).
\item \textbf{(Forward noncollinearity.)} The forward direct and latent finite slopes \(\gamma,\rho_1,\ldots,\rho_m\) are pairwise distinct.
\item \textbf{(Nonzero forward edge.)} The forward direct slope satisfies \(\gamma\ne 0\).
\item \textbf{(Observed law.)} The law \(P\) is the pushforward \(\mu\)-law of \((X,Y)\), i.e. \(P=\mu\circ (X,Y)^{-1}\).
\end{itemize}
\end{definitionv}

\cref{obj:def:forward-lvlingam-class} is a representation-level class: membership means that the law has at least one centered non-Gaussian source representation satisfying the right-arrow axis convention.

\begin{definitionv}{def:reverse-lvlingam-class}[Reverse class \(\mathcal P^{\mathrm{left}}_{m,K}\)]
Using the ambient law class \(\operatorname{Laws}(\mathbb R^2)\) defined in \cref{obj:def:forward-lvlingam-class}, let \(\mathcal P^{\mathrm{left}}_{m,2m+2}\) be the set of all probability laws \(P\in \operatorname{Laws}(\mathbb R^2)\) for which there exist a probability space \((\Omega,\mathcal A,\mu)\), latent source variables
\[
S_0,\ldots,S_{m+1}:\Omega\to\mathbb R,
\]
observed variables \(X,Y:\Omega\to\mathbb R\), a reverse direct slope \(\delta\in\mathbb R\), and reverse latent slopes \(\sigma_1,\ldots,\sigma_m\in\mathbb R\), such that:
\begin{itemize}
\item \textbf{(Source independence.)} The sources \(S_0,\ldots,S_{m+1}\) are mutually independent under \(\mu\).
\item \textbf{(Finite cumulant order.)} Each source belongs to \(L^{2m+2}(\mu)\).
\item \textbf{(Source non-Gaussianity.)} Each source law \(\mu\circ S_j^{-1}\) is non-Gaussian.
\item \textbf{(Centering.)} Each source is centered: \(\int_\Omega S_j\,d\mu=0\) for all \(j=0,\ldots,m+1\).
\item \textbf{(Reverse representation.)} The bivariate observed outcome vector \((X,Y)^\top\) admits the reverse \(Y\to X\) latent linear non-Gaussian representation with loading vectors \(v_j\), pointwise on \(\Omega\).
\item \textbf{(Reverse noncollinearity.)} The reverse direct and latent finite slopes \(\delta,\sigma_1,\ldots,\sigma_m\) are pairwise distinct.
\item \textbf{(Nonzero reverse edge.)} The reverse direct slope satisfies \(\delta\ne 0\).
\item \textbf{(Observed law.)} The law \(P\) is the pushforward of \(\mu\) by \(\omega\mapsto (X(\omega),Y(\omega))\).
\end{itemize}
\end{definitionv}

\cref{obj:def:reverse-lvlingam-class} gives the opposite representation class under the same source-count convention. The paper's separation statements characterize when cumulants from one parametrization have an empty fiber under the other cumulant map.

The identifying information is expressed through joint cumulants. For a law \(P\), \(\operatorname{Cum}\) denotes the joint cumulant functional, \leanref{sym:kappa_{r,a}(P)}{$\kappa_{r,a}(P)$} is the coordinate with \(r-a\) copies of \(X\) and \(a\) copies of \(Y\), and \leanref{sym:T_L(P)}{$T_L(P)$} collects all cumulants from orders \(2\) through \leanref{sym:L}{$L$}. The dimension of this vector is \leanref{sym:q_L}{$q_L$}. Classical references for cumulant notation and higher-order moment methods include \citet{Brillinger1969} and \citet{McCullagh1987}.

\begin{definitionv}{def:truncated-cumulant}[Truncated cumulant vector \(T_L(P)\)]
For the ambient law class of \cref{obj:def:forward-lvlingam-class}, write \(\operatorname{Cum}_P(Z_1,\ldots,Z_r)\) for the joint cumulant under \(P\), and define \(\operatorname{Cum}_r(Z)\) to be the joint cumulant of \(r\) copies of \(Z\). For \(P\in \operatorname{Laws}(\mathbb R^2)\), define
\[
T_L(P)=\bigl(\kappa_{r,a}(P):2\le r\le L,\ 0\le a\le r\bigr),
\]
where
\[
\kappa_{r,a}(P)
=
\operatorname{Cum}_P(\underbrace{X,\ldots,X}_{r-a\ \mathrm{copies}},
\underbrace{Y,\ldots,Y}_{a\ \mathrm{copies}}).
\]
The coordinate dimension is
\[
q_L=\frac{L(L+3)}{2}-2.
\]

For later cumulant maps, define the complex parameter spaces
\[
\Theta^{\mathrm{right}}_{m,L}
=\mathbb C^{,m+1+(m+2)(L-1)}
\]
with coordinates \((\gamma,\rho_1,\ldots,\rho_m,(c_{jr})_{0\le j\le m+1,,2\le r\le L})\), and
\[
\Theta^{\mathrm{left}}_{m,L}
=\mathbb C^{,m+1+(m+2)(L-1)}
\]
with coordinates \((\delta,\sigma_1,\ldots,\sigma_m,(d_{jr})_{0\le j\le m+1,,2\le r\le L})\).
\end{definitionv}

\cref{obj:def:truncated-cumulant} also introduces the complex affine parameter spaces \leanref{sym:Theta^right_{m,L}}{$\Theta^{\mathrm{right}}_{m,L}$} and \leanref{sym:Theta^left_{m,L}}{$\Theta^{\mathrm{left}}_{m,L}$}. These complexified coordinates, summarized by \leanref{S-2}{the complexified structural-parameter and cumulant-coordinate world}, are algebraic devices for studying polynomial images; real feasibility is imposed separately below.

By independence and multilinearity of cumulants, each source contributes its order-\(r\) cumulant weight times the appropriate monomial in its loading vector. We write \leanref{sym:c_{jr}}{$c_{jr}$} and \leanref{sym:d_{jr}}{$d_{jr}$} for these source-cumulant weights, and use the polynomial maps \leanref{sym:Phi^right_{m,L}}{$\Phi^{\mathrm{right}}_{m,L}$} and \leanref{sym:Phi^left_{m,L}}{$\Phi^{\mathrm{left}}_{m,L}$} to encode the resulting cumulant vectors.

\begin{definitionv}{def:forward-cumulant-map}[Forward cumulant map \(\Phi^{\mathrm{right}}_{m,L}\)]
Using the complex parameter space of \cref{obj:def:truncated-cumulant}, The forward cumulant map \(\Phi^{\mathrm{right}}_{m,L}:\Theta^{\mathrm{right}}_{m,L}\to\mathbb C^{q_L}\) is defined coordinatewise by
\[
\bigl[\Phi^{\mathrm{right}}_{m,L}(\gamma,\rho,c)\bigr]_{r,a}
=
\sum_{j=0}^{m+1}
c_{jr}\,(u_{j1})^{r-a}(u_{j2})^a,
\qquad
2\le r\le L,\quad 0\le a\le r,
\]
where
\[
u_0=(1,\gamma),\qquad
u_j=(1,\rho_j)\ \text{for }1\le j\le m,\qquad
u_{m+1}=(0,1).
\]
\end{definitionv}

The map in \cref{obj:def:forward-cumulant-map} is the cumulant analogue of a finite simultaneous binary-form decomposition. At this stage its source weights are free algebraic coordinates, with real-random-variable realizability imposed later through the feasible regions.

\begin{definitionv}{def:reverse-cumulant-map}[Reverse cumulant map \(\Phi^{\mathrm{left}}_{m,L}\)]
Using the complex parameter space of \cref{obj:def:truncated-cumulant}, The reverse cumulant map \(\Phi^{\mathrm{left}}_{m,L}:\Theta^{\mathrm{left}}_{m,L}\to\mathbb C^{q_L}\) is defined coordinatewise by
\[
\bigl[\Phi^{\mathrm{left}}_{m,L}(\delta,\sigma,d)\bigr]_{r,a}
=
\sum_{j=0}^{m+1}
d_{jr}\,(v_{j1})^{r-a}(v_{j2})^a,
\qquad
2\le r\le L,\quad 0\le a\le r,
\]
where
\[
v_0=(1,0),\qquad
v_j=(\sigma_j,1)\ \text{for }1\le j\le m,\qquad
v_{m+1}=(\delta,1).
\]
\end{definitionv}

\cref{obj:def:reverse-cumulant-map} mirrors \cref{obj:def:forward-cumulant-map} after exchanging the axis convention. This symmetry lets the paper characterize possible simultaneous membership in both parametrized images.

The algebraic images are taken after Zariski closure in the ambient complex affine cumulant space. We write \(\operatorname{ZariskiClosure}\) for this closure operation, and denote the two closed image varieties by \leanref{sym:C^right_{m,L}}{$C^{\mathrm{right}}_{m,L}$} and \leanref{sym:C^left_{m,L}}{$C^{\mathrm{left}}_{m,L}$}.

\begin{definitionv}{def:image-varieties}[Image varieties \(C^{\mathrm{right}}_{m,L}\) and \(C^{\mathrm{left}}_{m,L}\)]
For a subset of the ambient complex affine space, write \(\operatorname{ZariskiClosure}\) for its Zariski closure. Using the parameter spaces and maps of \cref{obj:def:forward-cumulant-map} and~\cref{obj:def:reverse-cumulant-map}, the forward and reverse image varieties are
\[
C^{\mathrm{right}}_{m,L}
=
\operatorname{ZariskiClosure}
\bigl(\Phi^{\mathrm{right}}_{m,L}(\Theta^{\mathrm{right}}_{m,L})\bigr),
\qquad
C^{\mathrm{left}}_{m,L}
=
\operatorname{ZariskiClosure}
\bigl(\Phi^{\mathrm{left}}_{m,L}(\Theta^{\mathrm{left}}_{m,L})\bigr),
\]
as subvarieties of \(\mathbb C^{q_L}\).
\end{definitionv}

These varieties record the algebraic restrictions imposed by a fixed number of source slots and a fixed arrow convention. Their use is standard in algebraic treatments of tensor and binary-form decompositions, where generic conclusions are naturally stated off proper algebraic subsets \citep{ComonMourrain1996,Landsberg2012}.

The next definition isolates the open parameter loci on which slopes are distinct, direct slopes are nonzero, and all retained source cumulant weights are nonzero. We denote these loci by \leanref{sym:Theta^{right,circ}_{m,L}}{$\Theta^{\mathrm{right},\circ}_{m,L}$} and \leanref{sym:Theta^{left,circ}_{m,L}}{$\Theta^{\mathrm{left},\circ}_{m,L}$}.

\begin{definitionv}{def:generic-parameter-loci}[Generic loci \(\Theta^{\mathrm{right},\circ}_{m,L}\) and \(\Theta^{\mathrm{left},\circ}_{m,L}\)]
For every \(L\ge2\), using the parameter spaces of \cref{obj:def:forward-cumulant-map} and~\cref{obj:def:reverse-cumulant-map}, the forward generic parameter locus is
\[
\Theta^{\mathrm{right},\circ}_{m,L}
=
\left\{
(\gamma,\rho,c)\in \Theta^{\mathrm{right}}_{m,L}:
\gamma
\prod_{i=1}^m(\gamma-\rho_i)
\prod_{1\le i<\ell\le m}(\rho_i-\rho_\ell)
\prod_{j=0}^{m+1}\prod_{r=2}^L c_{jr}
\ne 0
\right\}.
\]
The reverse generic parameter locus is
\[
\Theta^{\mathrm{left},\circ}_{m,L}
=
\left\{
(\delta,\sigma,d)\in \Theta^{\mathrm{left}}_{m,L}:
\delta
\prod_{i=1}^m(\delta-\sigma_i)
\prod_{1\le i<\ell\le m}(\sigma_i-\sigma_\ell)
\prod_{j=0}^{m+1}\prod_{r=2}^L d_{jr}
\ne 0
\right\}.
\]
\end{definitionv}

\cref{obj:def:generic-parameter-loci} strengthens the assumption-level genericity by also requiring nonzero retained source cumulants. This representation-level restriction selects the finite cumulant coordinates used by the generic theorem.

For a cumulant coordinate vector \leanref{sym:t}{$t$}, the relevant inverse images are the parameter fibers under the two polynomial maps. We use \leanref{sym:R^right_{m,L}(t)}{$R^{\mathrm{right}}_{m,L}(t)$} and \leanref{sym:R^left_{m,L}(t)}{$R^{\mathrm{left}}_{m,L}(t)$} for these fibers.

\begin{definitionv}{def:fiber-correspondences}
For a cumulant vector \(t\), the forward and reverse fibers are
\[
R^{\mathrm{right}}_{m,L}(t)
=
\left\{
\theta\in \Theta^{\mathrm{right}}_{m,L}:
\Phi^{\mathrm{right}}_{m,L}(\theta)=t
\right\},
\qquad
R^{\mathrm{left}}_{m,L}(t)
=
\left\{
\eta\in \Theta^{\mathrm{left}}_{m,L}:
\Phi^{\mathrm{left}}_{m,L}(\eta)=t
\right\}.
\]
\end{definitionv}

\cref{obj:def:fiber-correspondences} separates two questions: whether a cumulant vector lies in an arrow image, and how many parameter points produce it. The same-arrow fiber facts in the next section therefore have a different status from opposite-arrow exclusion.

Source labels in the middle block are exchangeable. The admissible relabelling group \leanref{sym:G_m}{$G_m$} acts by a permutation \leanref{sym:pi}{$\pi$} of the middle labels while fixing the two axis slots.

\begin{definitionv}{def:admissible-source-swaps}
Let \(G_m\) be the permutation group of the middle source labels \(\{1,\ldots,m\}\), acting while fixing labels \(0\) and \(m+1\). For \(\pi\in G_m\), the forward action on parameters is
\[
\pi\cdot(\gamma,(\rho_i)_{i=1}^m,(c_{jr})_{j,r})
=
(\gamma,(\rho_{\pi(i)})_{i=1}^m,(c'_{jr})_{j,r}),
\]
where
\[
c'_{ir}=c_{\pi(i),r}\quad(1\le i\le m),\qquad
c'_{0r}=c_{0r},\qquad
c'_{m+1,r}=c_{m+1,r}.
\]
The left action is the analogous action with \(\gamma,\rho,c\) replaced by \(\delta,\sigma,d\). Thus the source labels \(0\) and \(m+1\) are fixed, while the labels \(1,\ldots,m\) are permuted.
\end{definitionv}

\cref{obj:def:admissible-source-swaps} specifies the label symmetry quotiented in the paper. It fixes the direct-axis and fixed-axis source slots and permits permutations among the middle latent-source slots.

\subsection{The real-feasibility bridge}

The algebraic parameter spaces allow arbitrary complex source-cumulant weights. For econometric interpretation, we also need real parameter points whose source cumulants are realized by centered non-Gaussian real random variables through the retained order. These real feasible regions are denoted \leanref{sym:F^right_{m,L}}{$F^{\mathrm{right}}_{m,L}$} and \leanref{sym:F^left_{m,L}}{$F^{\mathrm{left}}_{m,L}$}. They bridge the complex theorem layer to real source models and retain the theorem's parameter-space interpretation.

\begin{definitionv}{def:real-feasible-regions}[Real feasible regions \(F^{\mathrm{right}}_{m,L}\) and \(F^{\mathrm{left}}_{m,L}\)]
Using the cumulant convention of \cref{obj:def:truncated-cumulant}, the forward real feasible region \(F^{\mathrm{right}}_{m,L}\) is the set of all real tuples \((\gamma,\rho,c)\in \mathbb R^{m+1}\times \mathbb R^{(m+2)(L-1)}\) such that \(\gamma\neq 0\), the slopes \(\rho_1,\ldots,\rho_m,\gamma\) are pairwise distinct, and, for every \(0\leq j\leq m+1\), there exists a centered non-Gaussian real random variable \(S_j\) with \(\mathbb E|S_j|^L<\infty\) satisfying
\[
\operatorname{Cum}_r(S_j)=c_{jr},
\qquad 2\leq r\leq L.
\]
The reverse real feasible region \(F^{\mathrm{left}}_{m,L}\) is defined analogously, with \((\delta,\sigma,d)\) in place of \((\gamma,\rho,c)\) and with \(\delta\neq 0\).
\end{definitionv}

\cref{obj:def:real-feasible-regions} identifies the finite real cumulant coordinates realized by actual centered non-Gaussian source laws.

We summarize the amount of cumulant information needed for generic real arrow separation by the information order \leanref{sym:K^star(m)}{$K^\star(m)$}. The arrow index \leanref{sym:b}{$b$} ranges over the two conventions, and \(\operatorname{opposite}(b)\) denotes the opposite-arrow involution.

\begin{definitionv}{def:information-order}[Information order \(K^\star(m)\)]
For \(b\in\{\mathrm{right},\mathrm{left}\}\), use the branch convention \(\Phi^b_{m,L}=\Phi^{\mathrm{right}}_{m,L}\) and \(F^b_{m,L}=F^{\mathrm{right}}_{m,L}\) when \(b=\mathrm{right}\), and \(\Phi^b_{m,L}=\Phi^{\mathrm{left}}_{m,L}\) and \(F^b_{m,L}=F^{\mathrm{left}}_{m,L}\) when \(b=\mathrm{left}\), using \cref{obj:def:forward-cumulant-map},~\cref{obj:def:reverse-cumulant-map}, and~\cref{obj:def:real-feasible-regions}. Define \(\operatorname{opposite}(\mathrm{right})=\mathrm{left}\) and \(\operatorname{opposite}(\mathrm{left})=\mathrm{right}\). Define \(K^\star(m)\) to be the minimum integer \(L\geq 2\), if such an integer exists, such that for each \(b\in\{\mathrm{right},\mathrm{left}\}\), outside a proper real algebraic subset of \(F^b_{m,L}\), the cumulant vector
\[
\Phi^b_{m,L}(\theta)
\]
has no representation in
\[
F^{\operatorname{opposite}(b)}_{m,L}.
\]
If no such \(L\) exists, set \(K^\star(m)=\infty\).
\end{definitionv}

\cref{obj:def:information-order} introduces notation for the generic separation threshold. The value \(K^\star(m)=2m+2\) serves as a conjectural benchmark, and the later information-order theorem establishes the sharper value in the relevant range.

Finally, the main theorem uses the classical apolar language for binary forms. For a finite projective support \(D\), the polynomial \(Q_D\) annihilates that support, and \(q(\partial)\) denotes the associated constant-coefficient differential operator. This notation is standard in binary-form decomposition and Waring-rank geometry \citep{ComonMourrain1996,Landsberg2012}.

\begin{definitionv}{def:apolar-notation}[Apolar operators \(Q_D\) and \(q(\partial)\)]
For a finite projective loading-direction support \(D\), define \(Q_D\) to be its squarefree homogeneous support-annihilator polynomial, with the normalization used by the cumulant-map coordinates. For a homogeneous binary form \(q(x,y)\), define \(q(\partial)\) to be the constant-coefficient differential operator obtained by substituting \(\partial_x,\partial_y\) for \(x,y\). For a real parameter set \(A\), define \(A_{\mathbb C}\) to be its coordinatewise complexification.
\end{definitionv}

\cref{obj:def:apolar-notation} supplies the notation used to recover unordered loading directions from truncated higher-order population cumulants.

\section{Main results}

The setup in \cref{obj:def:apolar-notation} turns the direction problem into a question about finite collections of binary forms. At order \(K=2m+2\), the cumulant coordinates determine a simultaneous decomposition with \(m+2\) loading directions. The first result states that, on a relatively Zariski-open dense set of parameters, the apolar equations recover the unordered loading support and rule out a full representation in the opposite arrow convention. This representation-level separation statement concerns the fibers of the polynomial cumulant maps in \cref{obj:def:forward-cumulant-map} and~\cref{obj:def:reverse-cumulant-map}.

\begin{theoremv}{thm:generic-apolar-arrow-recovery}
Let \(m\geq 1\), put \(n=m+2\) and \(K=2m+2=2n-2\). Use the support-annihilator and differential-operator conventions of \cref{obj:def:apolar-notation}.  There exist sets
\[
U_m^{\mathrm{right}}\subseteq \Theta^{\mathrm{right},\circ}_{m,K},
\qquad
U_m^{\mathrm{left}}\subseteq \Theta^{\mathrm{left},\circ}_{m,K}
\]
such that:

\begin{itemize}
\item \textbf{(Generic loci.)} Each of \(U_m^{\mathrm{right}}\) and \(U_m^{\mathrm{left}}\) is relatively Zariski open and relatively Zariski dense in the corresponding complex parameter space \(\Theta^{\mathrm{right}}_{m,K}\) or \(\Theta^{\mathrm{left}}_{m,K}\) of \cref{obj:def:truncated-cumulant}, and is contained in the corresponding generic parameter locus of \cref{obj:def:generic-parameter-loci}.

\item \textbf{(Open neighborhoods meeting feasibility.)} There are Euclidean-open sets \(O^{\mathrm{right}},O^{\mathrm{left}}\) of real parameter space such that
\[
O^{\mathrm{right}}\cap F^{\mathrm{right}}_{m,K}\neq\varnothing,
\qquad
O^{\mathrm{left}}\cap F^{\mathrm{left}}_{m,K}\neq\varnothing,
\]
where the feasible regions are those of \cref{obj:def:real-feasible-regions}; use the coordinatewise-complexification convention of \cref{obj:def:apolar-notation}. Then
\[
\left(O^{\mathrm{right}}\cap F^{\mathrm{right}}_{m,K}\right)_{\mathbb C}
 \subseteq U_m^{\mathrm{right}},
\qquad
\left(O^{\mathrm{left}}\cap F^{\mathrm{left}}_{m,K}\right)_{\mathbb C}
 \subseteq U_m^{\mathrm{left}}.
\]

\item \textbf{(Forward recovery and separation.)} For every \(\theta\in U_m^{\mathrm{right}}\),
\[
R^{\mathrm{left}}_{m,K}\!\left(\Phi^{\mathrm{right}}_{m,K}(\theta)\right)=\varnothing.
\]
Moreover, there is a nonzero squarefree polynomial \(Q\in\mathbb C[z]\) such that
\(Q(0)\neq0\) and the multiset of roots of \(Q\) is
\(\{\gamma,\rho_1,\ldots,\rho_m\}\). Every generic forward parameter \(\theta''\) satisfying
\[
\Phi^{\mathrm{right}}_{m,K}(\theta'')=
\Phi^{\mathrm{right}}_{m,K}(\theta)
\]
has the same multiset \(\{\gamma,\rho_1,\ldots,\rho_m\}\).  If
\[
f_r(x,y)=\sum_{a=0}^{r}\binom{r}{a}t_{r,a}x^{r-a}y^a,
\qquad
t=\Phi^{\mathrm{right}}_{m,K}(\theta),
\]
let \(D\) be the forward loading-direction support. Then
\[
Q_D\neq0,\qquad x\mid Q_D,\qquad y\nmid Q_D,
\]
and, for every homogeneous binary form \(q\) of degree \(n\),
\[
\bigl[q(\partial)f_{n+k}=0\ \text{for every }0\leq k\leq m\bigr]
\quad\Longleftrightarrow\quad
\exists c\in\mathbb C:\ q=cQ_D.
\]

\item \textbf{(Reverse recovery and separation.)} For every \(\eta\in U_m^{\mathrm{left}}\),
\[
R^{\mathrm{right}}_{m,K}\!\left(\Phi^{\mathrm{left}}_{m,K}(\eta)\right)=\varnothing.
\]
Moreover, there is a nonzero squarefree polynomial \(Q\in\mathbb C[z]\) such that
\(Q(0)\neq0\) and the multiset of roots of \(Q\) is
\(\{\delta,\sigma_1,\ldots,\sigma_m\}\). Every generic reverse parameter \(\eta''\) satisfying
\[
\Phi^{\mathrm{left}}_{m,K}(\eta'')=
\Phi^{\mathrm{left}}_{m,K}(\eta)
\]
has the same multiset \(\{\delta,\sigma_1,\ldots,\sigma_m\}\).  If
\[
f_r(x,y)=\sum_{a=0}^{r}\binom{r}{a}t_{r,a}x^{r-a}y^a,
\qquad
t=\Phi^{\mathrm{left}}_{m,K}(\eta),
\]
let \(D\) be the reverse loading-direction support. Then
\[
Q_D\neq0,\qquad y\mid Q_D,\qquad x\nmid Q_D,
\]
and, for every homogeneous binary form \(q\) of degree \(n\),
\[
\bigl[q(\partial)f_{n+k}=0\ \text{for every }0\leq k\leq m\bigr]
\quad\Longleftrightarrow\quad
\exists c\in\mathbb C:\ q=cQ_D.
\]
\end{itemize}
\end{theoremv}

The feasible-incidence clause in \cref{obj:thm:generic-apolar-arrow-recovery} supplies real feasible parameter points inside the open patches to which the complex-algebraic conclusion applies after coordinatewise complexification. Its topological content is intersection with the feasible regions \(F^{\mathrm{right}}_{m,K}\) and \(F^{\mathrm{left}}_{m,K}\), giving a precise real-parameter interpretation of the generic result.

The recovery statement is closely related to classical apolarity for binary forms and simultaneous tensor decompositions \citep{ComonMourrain1996,Landsberg2012,KoldaBader2009}. Here the distinctive identification point is that the recovered support contains one of the two fixed axes in a way that is incompatible with the opposite arrow on the generic loci. Thus, within the fixed source-count, axis-normalized representation class, the higher-order cumulants jointly recover loading directions and separate the two truncated-cumulant fibers.

The next result isolates the same-arrow consequences of \cref{obj:thm:generic-apolar-arrow-recovery} on its loci \(U_m^{\mathrm{right}}\) and \(U_m^{\mathrm{left}}\).

\begin{theoremv}{thm:generic-arrow-recovery-and-fiber-obstruction}[Generic Fiber Obstruction]
For every \(m\in\mathbb N\) with \(m\ge 1\), set \(K=2m+2\).  There exist complex loci
\(U^{\mathrm{right}}_m\subseteq \Theta^{\mathrm{right},\circ}_{m,K}\) and
\(U^{\mathrm{left}}_m\subseteq \Theta^{\mathrm{left},\circ}_{m,K}\), with the generic loci as in \cref{obj:def:generic-parameter-loci}, such that:
\begin{itemize}
\item \textbf{(Relative genericity.)}  Each of \(U^{\mathrm{right}}_m\) and \(U^{\mathrm{left}}_m\) is relatively Zariski open and relatively Zariski dense in the parameter spaces \(\Theta^{\mathrm{right}}_{m,K}\) and \(\Theta^{\mathrm{left}}_{m,K}\) of \cref{obj:def:forward-cumulant-map} and~\cref{obj:def:reverse-cumulant-map}, respectively.

\item \textbf{(Real incidence.)}  There is a Euclidean-open set of real forward parameters whose intersection with \(F^{\mathrm{right}}_{m,K}\) from \cref{obj:def:real-feasible-regions} is nonempty, and every parameter in this intersection has coordinatewise complexification contained in \(U^{\mathrm{right}}_m\).  Likewise, there is a Euclidean-open set of real reverse parameters whose intersection with \(F^{\mathrm{left}}_{m,K}\) is nonempty, and every parameter in this intersection has coordinatewise complexification contained in \(U^{\mathrm{left}}_m\).

\item \textbf{(Forward generic fibers.)}  For every \(\theta=(\gamma,\rho,c)\in U^{\mathrm{right}}_m\), every forward parameter in the forward fiber
\[
R^{\mathrm{right}}_{m,K}\!\left(\Phi^{\mathrm{right}}_{m,K}(\theta)\right)
\]
of \cref{obj:def:fiber-correspondences} and \cref{obj:def:forward-cumulant-map} has the same unordered loading-slope multiset
\[
\{\gamma,\rho_1,\ldots,\rho_m\}
\]
as \(\theta\).  The opposite-arrow fiber
\[
R^{\mathrm{left}}_{m,K}\!\left(\Phi^{\mathrm{right}}_{m,K}(\theta)\right)
\]
of \cref{obj:def:fiber-correspondences} (formed with the reverse map of \cref{obj:def:reverse-cumulant-map}) is empty.  For every latent index \(i=1,\ldots,m\), there is a generic forward parameter \(\theta'\in\Theta^{\mathrm{right},\circ}_{m,K}\) obtained by swapping the full pairs \((\gamma,(c_{0r})_{r=2}^K)\) and \((\rho_i,(c_{ir})_{r=2}^K)\), leaving every other coordinate fixed, such that
\[
\Phi^{\mathrm{right}}_{m,K}(\theta')
=
\Phi^{\mathrm{right}}_{m,K}(\theta),
\]
but \(\theta'\) is not in the admissible \(G_m\)-orbit of \(\theta\) from \cref{obj:def:admissible-source-swaps}.  If \(m\ge 2\), this forward fiber has exact relative Zariski dimension
\[
\frac{m(m-1)}{2}.
\]

\item \textbf{(Reverse generic fibers.)}  For every \(\eta=(\delta,\sigma,d)\in U^{\mathrm{left}}_m\), every reverse parameter in the reverse fiber
\[
R^{\mathrm{left}}_{m,K}\!\left(\Phi^{\mathrm{left}}_{m,K}(\eta)\right)
\]
has the same unordered loading-slope multiset
\[
\{\delta,\sigma_1,\ldots,\sigma_m\}
\]
as \(\eta\).  The opposite-arrow fiber
\[
R^{\mathrm{right}}_{m,K}\!\left(\Phi^{\mathrm{left}}_{m,K}(\eta)\right)
\]
is empty.  For every latent index \(i=1,\ldots,m\), there is a generic reverse parameter \(\eta'\in\Theta^{\mathrm{left},\circ}_{m,K}\) obtained by swapping the full pairs \((\delta,(d_{m+1,r})_{r=2}^K)\) and \((\sigma_i,(d_{ir})_{r=2}^K)\), leaving every other coordinate fixed, such that
\[
\Phi^{\mathrm{left}}_{m,K}(\eta')
=
\Phi^{\mathrm{left}}_{m,K}(\eta),
\]
but \(\eta'\) is not in the admissible \(G_m\)-orbit of \(\eta\) from \cref{obj:def:admissible-source-swaps}.  If \(m\ge 2\), this reverse fiber has exact relative Zariski dimension
\[
\frac{m(m-1)}{2}.
\]
\end{itemize}
\end{theoremv}

The non-orbit clause in \cref{obj:thm:generic-arrow-recovery-and-fiber-obstruction} is important for interpreting same-arrow identification. The admissible relabelling group fixes the axis slots, while the cumulant equations can still admit generic same-arrow parameter points obtained by exchanging a direct slot with a middle slot. Consequently, the theorem combines generic arrow separation with an exact witness of residual structural ambiguity within the chosen arrow convention.

We next define the exceptional compatibility locus. It consists of cumulant vectors that have a generic representation for one arrow and at least one full representation for the other. Its formulation in cumulant space makes the exceptional set a property of the truncated cumulant vector, independent of a particular chosen parametrization.

\begin{definitionv}{def:generic-full-fiber-compatibility}[Generic full-fiber compatibility \(E_m\)]
Set \(K=2m+2\). Define
\[
E_m
=
\left\{
t\in \mathbb C^{q_K}:
\bigl(R^{\mathrm{right}}_{m,K}(t)\cap \Theta^{\mathrm{right},\circ}_{m,K}\neq\varnothing
\ \text{and}\ 
R^{\mathrm{left}}_{m,K}(t)\neq\varnothing\bigr)
\ \text{or}\
\bigl(R^{\mathrm{left}}_{m,K}(t)\cap \Theta^{\mathrm{left},\circ}_{m,K}\neq\varnothing
\ \text{and}\ 
R^{\mathrm{right}}_{m,K}(t)\neq\varnothing\bigr)
\right\}.
\]
With \(\operatorname{ZariskiClosure}\) as defined in \cref{obj:def:image-varieties}, let
\[
\overline E_m=\operatorname{ZariskiClosure}(E_m).
\]
The corresponding exceptional preimages are
\[
H^{\mathrm{right}}_m
=
\Theta^{\mathrm{right},\circ}_{m,K}
\cap
\bigl(\Phi^{\mathrm{right}}_{m,K}\bigr)^{-1}(E_m),
\qquad
H^{\mathrm{left}}_m
=
\Theta^{\mathrm{left},\circ}_{m,K}
\cap
\bigl(\Phi^{\mathrm{left}}_{m,K}\bigr)^{-1}(E_m).
\]
\end{definitionv}

For the smallest source counts, the compatibility equations can be written explicitly. These displayed systems give finite incidence descriptions of the algebraic event in which the two arrow parametrizations agree on the retained cumulants.

\begin{definitionv}{def:worked-compatibility-instances}
For \(m=1\) and \(K=4\), \(A_1\) is the full solution set, in the displayed forward coordinates, reverse coordinates, and cumulant coordinates \(t\), of the \(12\)-equation incidence system
\[
t_{r,a}
=
c_{0r}\gamma^a+c_{1r}\rho^a+c_{2r}\mathbb 1_{\{a=r\}}
=
d_{0r}\mathbb 1_{\{a=0\}}+d_{1r}\sigma^{r-a}+d_{2r}\delta^{r-a},
\qquad
r=2,3,4,\quad 0\leq a\leq r,
\]
subject to the disjunction that the forward or reverse coordinates lie in the corresponding generic locus of \cref{obj:def:generic-parameter-loci}. For \(p\in A_1\), write \(p^{\mathrm{right}}\) for its forward parameter coordinates. Its common-axis subfamily is given by
\[
\rho=\sigma=0,\qquad
\delta\gamma=1,\qquad
d_{0r}=c_{1r},\qquad
d_{1r}=c_{2r},\qquad
d_{2r}=c_{0r}\gamma^r.
\]

For \(m=2\) and \(K=6\), \(A_2\) is the full solution set, in the displayed forward coordinates, reverse coordinates, and cumulant coordinates \(t\), of the \(25\)-equation incidence system
\[
t_{r,a}
=
c_{0r}\gamma^a+c_{1r}\rho_1^a+c_{2r}\rho_2^a+c_{3r}\mathbb 1_{\{a=r\}}
=
d_{0r}\mathbb 1_{\{a=0\}}+d_{1r}\sigma_1^{r-a}+d_{2r}\sigma_2^{r-a}+d_{3r}\delta^{r-a},
\qquad
r=2,\ldots,6,\quad 0\leq a\leq r,
\]
subject to the same generic-locus disjunction. For \(p\in A_2\), \(p^{\mathrm{right}}\) again denotes its forward parameter coordinates. Its common-axis subfamily is given by
\[
\rho_1=\sigma_1=0,\qquad
\delta\gamma=1,\qquad
\sigma_2\rho_2=1,\qquad
d_{0r}=c_{1r},\qquad
d_{1r}=c_{3r},\qquad
d_{2r}=c_{2r}\rho_2^r,\qquad
d_{3r}=c_{0r}\gamma^r.
\]
\end{definitionv}

The codimension calculation compares simultaneous decompositions after removing only the admissible middle-source label symmetry. This quotient-fiber comparison is the algebraic object behind the exceptional-locus statement.

\begin{definitionv}{def:separation-handle}[Quotient-fiber comparison]
The separation handle is the algebraic comparison of simultaneous binary-form decompositions
\[
\Phi^{\mathrm{right}}_{m,K}(\theta)=\Phi^{\mathrm{left}}_{m,K}(\eta)
\]
after quotienting each same-arrow fiber by the admissible source-label permutation group \(G_m\) of \cref{obj:def:admissible-source-swaps}.  The comparison uses the axis directions fixed in the two cumulant maps and the Jacobian ranks of the resulting quotient-fiber equations.
\end{definitionv}

The next theorem says that the compatibility locus is exceptional and has codimension exactly one in each arrow image variety at the apolar order. For \(m\ge2\), this single-equation geometry is substantially larger than a count based on the \(m\) middle sources might suggest.

\begin{theoremv}{thm:exceptional-locus-codimension-one}
For every \(m\geq 1\), at \(K=2m+2\), let \(\overline E_m\), \(H^{\mathrm{right}}_m\), and \(H^{\mathrm{left}}_m\) be as in \cref{obj:def:generic-full-fiber-compatibility}, and let \(C^{\mathrm{right}}_{m,K}\) and \(C^{\mathrm{left}}_{m,K}\) be as in \cref{obj:def:image-varieties}. Then:
\begin{itemize}
\item \textbf{(Codimension.)} The exceptional closure \(\overline E_m\) has codimension exactly \(1\) in each of \(C^{\mathrm{right}}_{m,K}\) and \(C^{\mathrm{left}}_{m,K}\). Moreover, if \(m\geq 2\), then \(\overline E_m\) does not have codimension \(m\) in either \(C^{\mathrm{right}}_{m,K}\) or \(C^{\mathrm{left}}_{m,K}\).

\item \textbf{(Forward preimage.)} For every \(\theta\in\Theta^{\mathrm{right}}_{m,K}\),
\[
\theta\in H^{\mathrm{right}}_m
\quad\Longleftrightarrow\quad
\theta\in\Theta^{\mathrm{right},\circ}_{m,K}
\ \text{and}\
R^{\mathrm{left}}_{m,K}\!\left(\Phi^{\mathrm{right}}_{m,K}(\theta)\right)\neq\varnothing .
\]

\item \textbf{(Reverse preimage.)} For every \(\eta\in\Theta^{\mathrm{left}}_{m,K}\),
\[
\eta\in H^{\mathrm{left}}_m
\quad\Longleftrightarrow\quad
\eta\in\Theta^{\mathrm{left},\circ}_{m,K}
\ \text{and}\
R^{\mathrm{right}}_{m,K}\!\left(\Phi^{\mathrm{left}}_{m,K}(\eta)\right)\neq\varnothing .
\]

\item \textbf{(Low-dimensional specializations.)} If \(m=1\), then
\[
E_1
=
\left\{
t:\ \exists p\text{ satisfying the system }A_1,\ 
\Phi^{\mathrm{right}}_{1,4}(p^{\mathrm{right}})=t
\right\}.
\]
If \(m=2\), then
\[
E_2
=
\left\{
t:\ \exists p\text{ satisfying the system }A_2,\ 
\Phi^{\mathrm{right}}_{2,6}(p^{\mathrm{right}})=t
\right\},
\]
where \(A_1,A_2\) are the compatibility instances of \cref{obj:def:worked-compatibility-instances}.
\end{itemize}
\end{theoremv}

\cref{obj:thm:exceptional-locus-codimension-one} gives the proper interpretation of the exceptional set. Generic separation in \cref{obj:thm:generic-apolar-arrow-recovery} is compatible with a codimension-one family of cumulant vectors that admit both arrow conventions. The statement also identifies the preimages \(H^{\mathrm{right}}_m\) and \(H^{\mathrm{left}}_m\) exactly as generic same-arrow parameters whose cumulant vector retains a full opposite-arrow fiber.

The preceding results use order \(2m+2\), the natural apolar order for recovering \(m+2\) loading directions. The real information-order question asks whether this order is minimal for generic real arrow separation. The following handle records the relevant qualification: through order \(2m+1\), real twin parameters with matching truncated cumulants exist.

\begin{definitionv}{def:real-twin-construction-handle}
The real twin construction handle is the proposition that there exist real forward and reverse parameters \(\theta\) and \(\eta\), each realized through order
\[
K_-=2m+1,
\]
by centered, non-Gaussian, compactly supported real sources
such that their truncated cumulant maps agree through order \(K_-\):
\[
\Phi^{\mathrm{right}}_{m,K_-}(\theta)=\Phi^{\mathrm{left}}_{m,K_-}(\eta).
\]
This definition records only the existence proposition; it does not include a construction or proof.
\end{definitionv}

For \(m\ge3\), however, the generic real separation threshold improves by one order relative to the apolar support-recovery order. The statement establishes generic separation over the real feasible regions of \cref{obj:def:real-feasible-regions}, while full loading-support recovery remains the target of the order-\(2m+2\) apolar theorem.

\begin{theoremv}{thm:improved-real-information-order}[Improved Information Order]
The theorem-local symbols \(E^{\mathrm{right}}\) and \(E^{\mathrm{left}}\) below are real parameter-space exceptional subsets and are unrelated to the cumulant-space compatibility locus \(E_m\). For every \(m\in\mathbb N\) with \(3\le m\), at truncation order \(L=2m+1\):
\begin{itemize}
\item \textbf{(Forward generic separation.)} There is a proper real algebraic subset of the real forward parameter space whose complement meets \(F^{\mathrm{right}}_{m,L}\); choose one such subset and denote it locally in this theorem by \(E^{\mathrm{right}}\). For every \(\theta\in F^{\mathrm{right}}_{m,L}\setminus E^{\mathrm{right}}\) there is no \(\eta\in F^{\mathrm{left}}_{m,L}\) such that
\[
\Phi^{\mathrm{right}}_{m,L}(\theta)=\Phi^{\mathrm{left}}_{m,L}(\eta).
\]
\item \textbf{(Reverse generic separation.)} There is a proper real algebraic subset of the real reverse parameter space whose complement meets \(F^{\mathrm{left}}_{m,L}\); choose one such subset and denote it locally in this theorem by \(E^{\mathrm{left}}\). For every \(\eta\in F^{\mathrm{left}}_{m,L}\setminus E^{\mathrm{left}}\) there is no \(\theta\in F^{\mathrm{right}}_{m,L}\) such that
\[
\Phi^{\mathrm{left}}_{m,L}(\eta)=\Phi^{\mathrm{right}}_{m,L}(\theta).
\]
\end{itemize}
Consequently, for the information order \(K^\star(m)\) of \cref{obj:def:information-order},
\[
K^\star(m)\le 2m+1
\qquad\text{and}\qquad
K^\star(m)\ne 2m+2 .
\]
\end{theoremv}

\cref{obj:thm:improved-real-information-order} establishes \(K^\star(m)\le 2m+1\) for \(m\ge3\) under the generic real separation criterion in \cref{obj:def:information-order}. Taken together, \cref{obj:thm:generic-apolar-arrow-recovery}--\cref{obj:thm:improved-real-information-order} separate three issues: support recovery at the apolar order, same-arrow fiber nonuniqueness within a fixed axis convention, and the lower real truncation order sufficient for generic exclusion of the opposite arrow.

\section{Discussion and extensions}

\subsection{Interpretation}

The results in the preceding section identify a representation-level exclusion restriction in bivariate latent linear non-Gaussian models. \Cref{obj:thm:generic-apolar-arrow-recovery} excludes the opposite axis-normalized cumulant-map fiber. Thus, if a structural model \leanref{sym:M}{$M$} induces observational law \leanref{sym:P_M}{$P_M$}, the theorems characterize the truncated cumulants of \(P_M\) under the fixed source-count and axis-normalization conventions of \cref{obj:def:forward-lvlingam-class} and~\cref{obj:def:reverse-lvlingam-class}.

This distinction matters for econometric interpretation. In a conventional nonparametric identification problem, one asks whether the observational law determines a structural object over a specified model class. Here the maintained class is more structured: independent non-Gaussian source slots, a fixed and externally specified number of slots, and two axis-normalized loading conventions. A natural user is therefore a researcher whose institutional or factor model fixes these inputs and who wants to know whether the competing recursive specification is compatible with the population cumulants. Within that class, higher-order cumulants impose restrictions beyond covariance information and can generically exclude the opposite truncated-cumulant fiber. Related recent work on algebraic and semialgebraic identification likewise distinguishes generic algebraic separation from full distributional identification \citep{MestersZwiernik2022,Virolainen2024}.

The codimension-one exceptional locus in \cref{obj:thm:exceptional-locus-codimension-one} is a compatibility set for truncated cumulant vectors. Generic separation holds off this Zariski-closed set, and the theorem gives it exactly codimension one inside each arrow image variety at the apolar order. Econometrically, opposite-arrow compatibility is therefore a genuine lower-dimensional restriction governed by a single algebraic dimension even when \(m\ge2\). The size of this set motivates explicit reporting in applications where the data-generating process may be close to a compatibility surface.

The same-arrow fiber statements in \cref{obj:thm:generic-arrow-recovery-and-fiber-obstruction} sharpen the interpretation of the information recovered. The unordered loading-slope support is recovered on the stated generic loci, while the chosen arrow retains parametrization ambiguity beyond the admissible middle-source relabellings of \cref{obj:def:admissible-source-swaps}. In particular, exchanging a direct slot with a middle slot can preserve the cumulant vector while falling outside the admissible \(G_m\)-orbit. Opposite-fiber exclusion is therefore compatible with substantial within-convention nonidentification.

The information-order result has a separate role. The apolar theorem uses order \(2m+2\) to recover the unordered loading support through the binary-form annihilator equations. \Cref{obj:thm:improved-real-information-order} shows that, for \(m\ge3\), generic real exclusion of the opposite arrow can already occur at order \(2m+1\). The two thresholds therefore align with distinct targets: generic real arrow separation at \(2m+1\) and loading-support recovery at \(2m+2\).

\subsection{Limitations and future work}

These qualifications also delimit the connection to empirical practice. The paper provides population-level algebraic identification statements for truncated cumulants. Its immediate user is a theorist or model-class analyst who treats \(m\) and the axis normalization as maintained inputs and needs to know whether the population cumulant map separates the two representations. It does not supply a finite-sample estimator, a sampling distribution, a model-selection criterion, or a numerical decision rule. Any applied implementation would have to confront cumulant estimation error, weak higher-order signal, near-exceptional parameter values, and uncertainty about the number of source slots. Recent work on latent non-Gaussian identification and causal discovery develops complementary estimation and testing perspectives, but those problems are not solved by the generic separation theorem alone \citep{XieHuangChenEtAl2023,MorinishiShimizu2025,ChenGuPengEtAl2025}.

The main contribution is algebraic econometric identification for a fixed, axis-normalized latent-source representation class. Higher-order cumulants generically contain enough information to exclude the opposite convention's cumulant-map fiber inside that class, the exceptional compatibility set is exactly codimension one at the apolar order, and lower real truncation can suffice for generic opposite-fiber exclusion when \(m\ge3\). These statements establish the population information available before estimation and provide a foundation for future finite-sample procedures.

\section*{Appendices}

\section{Appendix with proofs and auxiliary lemmas}

This appendix collects the auxiliary invariance statement and the external algebraic background used to organize the results in \cref{obj:thm:generic-apolar-arrow-recovery} through \cref{obj:thm:improved-real-information-order}. The first point is purely structural: the middle latent-source slots are exchangeable, while the two axis slots retain their arrow-specific roles. The admissible relabellings of \cref{obj:def:admissible-source-swaps} therefore preserve the arrow convention and the induced cumulant vector.

\begin{lemmav}{lem:admissible-swaps-preserve-direction}[Admissible Swaps Preserve Direction]
Let \(m\ge 1\), let \(2\le L\le 2m+2\), and let \(\pi\in G_m\) be as in \cref{obj:def:admissible-source-swaps}. Let \(\Theta^{\mathrm{right}}_{m,L}\) and \(\Theta^{\mathrm{left}}_{m,L}\) be the parameter spaces of \cref{obj:def:forward-cumulant-map} and~\cref{obj:def:reverse-cumulant-map}. Let
\(\theta^{\mathbb C}\in\Theta^{\mathrm{right}}_{m,L}\) and
\(\eta^{\mathbb C}\in\Theta^{\mathrm{left}}_{m,L}\) be complex parameters, let
\(\theta\in F^{\mathrm{right}}_{m,L}\) and
\(\eta\in F^{\mathrm{left}}_{m,L}\) be real parameters when feasibility is assumed, and let
\(P_{\mathrm{right}},P_{\mathrm{left}}\in\operatorname{Laws}(\mathbb R^2)\), with the law and cumulant conventions of \cref{obj:def:truncated-cumulant}. Then:

\begin{itemize}
\item \textbf{(Complex cumulant maps.)} The admissible source relabeling of
\cref{obj:def:admissible-source-swaps} leaves both complex cumulant maps unchanged:
\[
\Phi^{\mathrm{right}}_{m,L}(\pi\theta^{\mathbb C})
  =\Phi^{\mathrm{right}}_{m,L}(\theta^{\mathbb C}),
\qquad
\Phi^{\mathrm{left}}_{m,L}(\pi\eta^{\mathbb C})
  =\Phi^{\mathrm{left}}_{m,L}(\eta^{\mathbb C}).
\]

\item \textbf{(Real feasible regions.)} If \(\theta\in F^{\mathrm{right}}_{m,L}\), then
\(\pi\theta\in F^{\mathrm{right}}_{m,L}\); and if
\(\eta\in F^{\mathrm{left}}_{m,L}\), then
\(\pi\eta\in F^{\mathrm{left}}_{m,L}\), with feasibility understood as in
\cref{obj:def:real-feasible-regions}.

\item \textbf{(Forward witnesses.)} For any probability space satisfying the usual
measurability and integrability conditions, any centered independent non-Gaussian
sources \(S_0,\ldots,S_{m+1}\) satisfying
\cref{obj:ass:independent-sources},
\cref{obj:ass:source-nongaussianity}, and
\cref{obj:ass:finite-cumulants}, and any observed pair \((X,Y)\) with
law \(P_{\mathrm{right}}\), suppose \(\theta\in F^{\mathrm{right}}_{m,L}\), the
forward axis model of \cref{obj:ass:forward-axis-model} holds with
parameter \(\theta\), the conditions of
\cref{obj:ass:forward-noncollinearity} and
\cref{obj:ass:forward-nonzero-edge} hold, and
\[
\operatorname{Cum}_r(S_j)=c_{jr}
\qquad\text{for every }j\text{ and every }2\le r\le L .
\]
After relabeling the middle sources by \(\pi\), the relabeled sources remain
independent, centered, non-Gaussian, and finite through order \(2m+2\); they satisfy
the same forward axis model with parameter \(\pi\theta\), the forward
noncollinearity and nonzero-edge conditions still hold, and their cumulants satisfy
\[
\operatorname{Cum}_r(S_{\pi(j)})=(\pi c)_{jr}
\qquad\text{for every }j\text{ and every }2\le r\le L .
\]
Consequently \(P_{\mathrm{right}}\in\mathcal P^{\mathrm{right}}_{m,2m+2}\) as in
\cref{obj:def:forward-lvlingam-class}.

\item \textbf{(Reverse witnesses.)} The analogous reverse statement holds: for any
reverse representation of \(P_{\mathrm{left}}\) satisfying
\cref{obj:ass:independent-sources},
\cref{obj:ass:source-nongaussianity},
\cref{obj:ass:finite-cumulants},
\cref{obj:ass:reverse-axis-model},
\cref{obj:ass:reverse-noncollinearity}, and
\cref{obj:ass:reverse-nonzero-edge}, with
\(\eta\in F^{\mathrm{left}}_{m,L}\) and
\[
\operatorname{Cum}_r(S_j)=d_{jr}
\qquad\text{for every }j\text{ and every }2\le r\le L ,
\]
the relabeled sources remain valid reverse witnesses with parameter \(\pi\eta\),
their cumulants satisfy
\[
\operatorname{Cum}_r(S_{\pi(j)})=(\pi d)_{jr}
\qquad\text{for every }j\text{ and every }2\le r\le L ,
\]
and \(P_{\mathrm{left}}\in\mathcal P^{\mathrm{left}}_{m,2m+2}\) as in
\cref{obj:def:reverse-lvlingam-class}.

\item \textbf{(Direction of structural models.)} With \(P_M\) and \(D(M)\) as defined in \cref{obj:def:separated-model-domain}, if a structural model \(M\) has
\(D(M)=X\to Y\) and parameter \(\theta\), then there is a structural model \(M'\)
with \(P_{M'}=P_M\), \(D(M')=X\to Y\), and parameter \(\pi\theta\).  If \(M\) has
\(D(M)=Y\to X\) and parameter \(\eta\), then there is a structural model \(M'\)
with \(P_{M'}=P_M\), \(D(M')=Y\to X\), and parameter \(\pi\eta\).

\item \textbf{(Arrow-tagged orbits.)} On complex parameters, the relabeled source
weights are exactly the permuted weights:
\[
(\pi c^{\mathbb C})_{jr}=c^{\mathbb C}_{\pi(j),r},
\qquad
(\pi d^{\mathbb C})_{jr}=d^{\mathbb C}_{\pi(j),r}.
\]
For the boundary source labels \(j=0\) and \(j=m+1\), the forward loading vectors
are fixed by the admissible relabeling.  Moreover every admissible relabeling
preserves the arrow tag: right-tagged parameters remain right-tagged and
left-tagged parameters remain left-tagged.  Hence the right-tagged \(G_m\)-orbit of
\(\theta^{\mathbb C}\) is disjoint from the left-tagged \(G_m\)-orbit of
\(\eta^{\mathbb C}\).

\item \textbf{(All admissible relabelings.)} For every \(\sigma\in G_m\),
\[
\Phi^{\mathrm{right}}_{m,L}(\sigma\theta^{\mathbb C})
  =\Phi^{\mathrm{right}}_{m,L}(\theta^{\mathbb C}),
\qquad
\Phi^{\mathrm{left}}_{m,L}(\sigma\eta^{\mathbb C})
  =\Phi^{\mathrm{left}}_{m,L}(\eta^{\mathbb C}).
\]
\end{itemize}
\end{lemmav}

\cref{obj:lem:admissible-swaps-preserve-direction} is the invariance fact needed whenever a same-arrow fiber is compared modulo source labels. It justifies quotienting by the middle-source permutation group in \cref{obj:def:separation-handle}, while keeping the direct-axis and fixed-axis slots attached to their arrow conventions. The non-orbit clauses in \cref{obj:thm:generic-arrow-recovery-and-fiber-obstruction} are therefore substantive: a direct slot can agree observationally with a middle slot in the same-arrow fiber and still lie beyond the admissible relabelling orbit.

\begin{remarkv}{def:real-atlas-handle}
The real atlas handle is the following external classical interface.  For every finite family of real polynomials and every variable order, there exists an adapted cylindrical algebraic decomposition that is sign-invariant for the generated projection family, is obtained by recursive section and sector lifting, decides polynomial sign conditions cellwise, and supports Tarski--Seidenberg semialgebraic projection.  This is the classical existence-and-structure interface of \citet[Section~2.3 and Theorem~2.2.1]{BochnakCosteRoy1998} and \citet[Chapter~5]{BasuPollackRoy2006}; it is paper-agnostic and carries no algorithmic, computability, payload, or complexity claim.  It is distinct from the exact rational Gröbner--CAD interface.
\end{remarkv}

\begin{remarkv}{def:effective-rational-groebner-cad-interface}[Effective rational Gröbner--CAD interface]
This paper assumes the following cited external interface. There exist natural constants \(a,b\) with \(a\ge 1\) and fixed uniform exact symbolic machines. For finite supplied polynomial presentations over \(\mathbb Q\) and \(\mathbb Q(i)\), supplied admissible monomial and variable orders, and supplied retained coordinate blocks, these machines terminate with certified Gröbner bases and normal forms, elimination ideals, ideal intersections, saturations, the rational complex elimination-closure theorem, and a sign-invariant real cylindrical algebraic decomposition with exact real-algebraic root and sign data, recursively lifted section-and-sector cells, quantifier elimination, cellwise truth and retention for polynomial equalities and weak or strict inequalities, and finite certified traces.

If \(D\ge 1\), every supplied job has ambient dimension at most \(N\) and supplied degree envelope at most \(D\), and the total supplied source-polynomial count is at most \(s\), one combined certificate bounds the sum of a finite batch of rational algebra traces, a finite batch of Gaussian-rational algebra traces, and the rational CAD and quantifier-elimination trace by
\[
\bigl(\max\{2,sD\}\bigr)^{2^{aN+b}}.
\]
The same fixed machines and constants also certify a dependent pipeline when the caller first supplies a jointly presentable cross-source coordinate relation. Two supplied saturated-elimination jobs finish, injective renamings identify their coordinates exactly according to that predeclared relation, their renamed outputs feed an ideal-intersection job, and an injective renaming sends the resulting intersection basis into a dependent CAD job. This certificate additionally requires the positive envelope \(D\) to bound the actual saturated-elimination bases, the intersection job and its actual basis, and the CAD input. This paper-agnostic interface assumes no LiNGAM map, moment certificate, exceptional locus, fiber, atlas cell, label, or conclusion.
\end{remarkv}

\section{Appendix: verification and reproducibility scope}

\begin{definitionv}{def:global-feasible-fiber-decision}[Order-\(2m+2\) feasible-fiber predicate \(\operatorname{FeasFiber}^{b}_{m,2m+2}(t)\)]
Set \(K=2m+2\). For \(b\in\{\mathrm{right},\mathrm{left}\}\), with the branch convention of \cref{obj:def:information-order}, define the order-\(K\) predicate \(\operatorname{FeasFiber}^{b}_{m,2m+2}(t)\) to be the atomic-certificate predicate asserting that there exist real \(b\)-loading and source-cumulant coordinates \(\lambda\), and, for each \(0\leq j\leq m+1\), witnesses
\[
(w_{jh},z_{jh})_{h=1}^{m+2},
\]
such that
\[
\Phi^b_{m,K}(\lambda)=t,
\]
the direct slope for \(b\) is nonzero, all finite loading slopes for \(b\) are pairwise distinct, and the following moment conditions hold. Write \(k_{jr}\) for the source cumulants contained in \(\lambda\), let \(B_r\) denote the complete exponential Bell polynomial (the cumulant-to-moment polynomial), and set
\[
\mu_{jr}=B_r(0,k_{j2},\ldots,k_{jr}).
\]
Then
\[
\sum_{h=1}^{m+2} w_{jh}=1,\qquad
w_{jh}\geq 0,\qquad
\sum_{h=1}^{m+2} w_{jh}z_{jh}^r=\mu_{jr}\quad(1\leq r\leq K),
\qquad
\mu_{j2}>0.
\]
Separately, the formal object records an unproved proposition asserting the existence of a sign-invariant cylindrical algebraic decomposition procedure, in variable order \(t\), then \(\lambda\), then the witnesses, that decides this displayed predicate for the fixed axis \(b\). The predicate definition itself does not assert equivalence with nonemptiness of the corresponding real-source feasible fiber.
\end{definitionv}

\begin{definitionv}{def:direction-selector}[Direction selector \(S_m\)]
Using the order-\(2m+2\) feasible-fiber predicates of \cref{obj:def:global-feasible-fiber-decision}, define the partial direction selector
\[
S_m(t)=
\begin{cases}
X\to Y,&\operatorname{FeasFiber}^{\mathrm{right}}_{m,2m+2}(t)\text{ holds and }\operatorname{FeasFiber}^{\mathrm{left}}_{m,2m+2}(t)\text{ does not},\\
Y\to X,&\operatorname{FeasFiber}^{\mathrm{left}}_{m,2m+2}(t)\text{ holds and }\operatorname{FeasFiber}^{\mathrm{right}}_{m,2m+2}(t)\text{ does not},\\
\text{undefined},&\text{otherwise}.
\end{cases}
\]
Separately, the formal object records the unproved proposition that a globally finite semialgebraic function on \(\mathbb R^{q_{2m+2}}\) agrees with \(S_m\) wherever exactly one feasibility predicate holds, together with Boolean decision functions for the forward and reverse maps. This definition supplies neither a proof of that proposition nor an evaluable decision procedure.
\end{definitionv}

\begin{definitionv}{def:separated-model-domain}[Separated model domain \(M^{\mathrm{sep}}_{m,K}\)]
For a structural model \(M\), write \(P_M\) for its induced observational law and \(D(M)\in\{X\to Y,Y\to X\}\) for its arrow. Define \(M^{\mathrm{sep}}_{m,K}\) to be the class of structural models \(M\) such that either \(M\) has a forward representation with parameter \(\theta\in F^{\mathrm{right}}_{m,K}\) and
\[
R^{\mathrm{left}}_{m,K}\bigl(T_K(P_M)\bigr)\cap F^{\mathrm{left}}_{m,K}=\varnothing,
\]
or \(M\) has a reverse representation with parameter \(\eta\in F^{\mathrm{left}}_{m,K}\) and
\[
R^{\mathrm{right}}_{m,K}\bigl(T_K(P_M)\bigr)\cap F^{\mathrm{right}}_{m,K}=\varnothing.
\]
\end{definitionv}

\begin{remarkv}{thm:exact-real-exceptional-atlas}[Exact Real Exceptional Atlas]
One worthwhile direction for future research is to determine whether \emph{For every fixed m, there is a finite computable semialgebraic cylindrical atlas of the real exceptional closure such that, on each base cell, the complete feasible forward and reverse fibers are finite unions of recursively cylindrical section and sector cells and their nonemptiness labels are constant. The atlas includes singular and nongeneric boundary branches and decides exactly when both real arrows are feasible.}. Doing so requires the remaining paper-specific exact-atlas constructor would require substantial CAD/elimination substrate, while the node is secondary and has no delivered theorem consumer, which lies beyond the present development.
\end{remarkv}

The formal claim in this paper is reproducible at the source revision used to build the manuscript. The repository is \url{https://github.com/Jiyuan-Tan/CausalForge}; the pinned revision is
\[
\texttt{2c9214dff2d737c1d131f6056916aad82bdc7e4b},
\]
and the CausalSmith project uses Lean \texttt{v4.29.0-rc3}. The matching
machine-readable contract is
\path{CausalSmith/doc/presentation/eid_lingam_direction_min_order_v1_truncated_cumulant_minimality/verification_contract.json};
the declaration table below refers to that same revision. From the repository
root, the checked project is built with
\[
\texttt{lake -d CausalSmith build}.
\]
The repository's \texttt{CausalSmith/lean-toolchain}, \texttt{CausalSmith/lakefile.toml}, and dependency manifests fix the toolchain and local Causalean dependency used by that build.

\begin{table}[ht]
\centering
\scriptsize
\caption{Main paper results and checked Lean declarations}
\begin{tabular}{p{0.30\linewidth}p{0.62\linewidth}}
\hline
Paper result & Lean declaration \\
\hline
Generic apolar arrow recovery &
\texttt{CausalSmith.ExactID.EID\_LingamDirectionMinOrderV1.generic\_apolar\_arrow\_recovery} \\
Same-arrow fiber consequences &
\texttt{CausalSmith.ExactID.EID\_LingamDirectionMinOrderV1.genericArrowRecoveryAndFiberObstruction} \\
Exceptional locus of codimension one &
\texttt{CausalSmith.ExactID.EID\_LingamDirectionMinOrderV1.exceptionalLocusCodimensionOne} \\
Improved real information order &
\texttt{CausalSmith.ExactID.EID\_LingamDirectionMinOrderV1.improvedRealInformationOrder} \\
Admissible-swap invariance &
\texttt{CausalSmith.ExactID.EID\_LingamDirectionMinOrderV1.admissibleSwaps\_preserve\_direction} \\
\hline
\end{tabular}
\end{table}

The generated \texttt{verification\_contract.json} in the paper bundle gives the full declaration-to-paper mapping, source file, statement match, and sorry-free status for every displayed formal object. The checked scope covers the frozen definitions and theorems used for the cumulant maps, generic separation, same-arrow fiber behavior, exceptional-locus dimension, and information-order result. Classical library facts and explicitly stated external algebraic inputs remain dependencies, while the prose supplies the surrounding interpretation.

\paragraph{Limitations.}
The formalization does not certify an estimator, sampling theory, selection of \(m\), or an exact decision procedure for all real feasible fibers, including singular and nongeneric boundary cases. Those tasks are outside the delivered contribution.

\section{Proofs of the main results}\label{sec:deferred-proofs}

\begin{proof}[Proof of \cref{obj:thm:generic-arrow-recovery-and-fiber-obstruction}]
Fix \(m\ge 1\) and put \(K=2m+2\).

\begin{enumerate}
\item \cref{obj:thm:generic-apolar-arrow-recovery} gives loci
\(U^{\mathrm{right}}_m\) and \(U^{\mathrm{left}}_m\), each contained in the corresponding generic locus of \cref{obj:def:generic-parameter-loci}, and supplies their relative Zariski openness, relative Zariski density, and the asserted real-incidence open sets.  For each forward point it also supplies an empty opposite-arrow fiber, a support-annihilator kernel identity for the divided-power forms \(f_r\), and a nonzero univariate polynomial whose roots are exactly the multiset
\(\{\gamma,\rho_1,\ldots,\rho_m\}\) and whose value at \(0\) is nonzero.  The reverse statement is supplied with the same data after interchanging the two arrows.

\item Let \(\theta=(\gamma,\rho,c)\in U^{\mathrm{right}}_m\).  If
\(\theta'\in R^{\mathrm{right}}_{m,K}(\Phi^{\mathrm{right}}_{m,K}(\theta))\), then the definition of the fiber gives equality of all cumulant coordinates in the active range \(2\le r\le K\), \(a\le r\); outside this range both polynomial maps are zero by definition, so
\[
\Phi^{\mathrm{right}}_{m,K}(\theta')
=
\Phi^{\mathrm{right}}_{m,K}(\theta).
\]
The support-annihilator of the loading directions of \(\theta'\) lies in the same contraction kernel.  By the kernel identity furnished in Step 1, it is a nonzero scalar multiple of the support-annihilator of \(\theta\).  Dehomogenizing along the forward finite-slope chart therefore gives the same root multiset, hence
\[
\{\gamma',\rho'_1,\ldots,\rho'_m\}
=
\{\gamma,\rho_1,\ldots,\rho_m\}.
\]
Thus every forward point in the forward fiber has the same unordered loading-slope multiset as \(\theta\).

\item The opposite-arrow conclusion for this \(\theta\) is exactly the empty reverse-fiber conclusion supplied in Step 1:
\[
R^{\mathrm{left}}_{m,K}(\Phi^{\mathrm{right}}_{m,K}(\theta))=\varnothing .
\]
Moreover, because the polynomial from Step 1 has roots \(\{\gamma,\rho_1,\ldots,\rho_m\}\) and is nonzero at \(0\), none of the latent slopes \(\rho_i\) is zero.

\item Fix \(i\in\{1,\ldots,m\}\), and define \(\theta'\) by swapping the full direct-source pair \((\gamma,(c_{0r})_{r=2}^K)\) with the full \(i\)-th latent pair \((\rho_i,(c_{ir})_{r=2}^K)\), leaving every other coordinate fixed.  Since \(\theta\) is generic and \(\rho_i\ne0\), the swapped parameter remains in \(\Theta^{\mathrm{right},\circ}_{m,K}\).

\item This \(\theta'\) is, by construction, the forward direct-latent swap required in the statement.

\item The forward cumulant map is unchanged under this full-pair swap: the defining source sum is merely reindexed by the transposition of the direct source label and the \(i\)-th latent source label.  Hence
\[
\Phi^{\mathrm{right}}_{m,K}(\theta')
=
\Phi^{\mathrm{right}}_{m,K}(\theta).
\]

\item The swapped parameter is not in the admissible \(G_m\)-orbit of \(\theta\).  Indeed, admissible relabelings only permute the middle latent source labels and keep the direct coordinate fixed, whereas the constructed swap changes the direct coordinate from \(\gamma\) to \(\rho_i\); genericity gives \(\gamma\ne\rho_i\).

\item If \(m\ge2\), apply the same-arrow fiber dimension statement to the generic point \(\theta\), using the multiset recovery proved in Step 2.  Once the loading multiset is recovered, the fiber is controlled by fixed-loading weight fibers; their independent low-order kernels have total dimension
\[
\sum_{r=2}^{m}(m+1-r)=\frac{m(m-1)}2,
\]
and every irreducible chain in the full fiber is contained in one of these fixed-loading components.  Therefore the forward fiber has exact relative Zariski dimension \(\frac{m(m-1)}2\).

\item Now let \(\eta=(\delta,\sigma,d)\in U^{\mathrm{left}}_m\).  If
\(\eta'\in R^{\mathrm{left}}_{m,K}(\Phi^{\mathrm{left}}_{m,K}(\eta))\), the same fiber-coordinate bookkeeping gives
\[
\Phi^{\mathrm{left}}_{m,K}(\eta')
=
\Phi^{\mathrm{left}}_{m,K}(\eta).
\]
The reverse support-annihilator of \(\eta'\) is then in the common reverse contraction kernel, hence is a nonzero scalar multiple of the reverse support-annihilator of \(\eta\).  Dehomogenizing in the reverse finite-slope chart gives
\[
\{\delta',\sigma'_1,\ldots,\sigma'_m\}
=
\{\delta,\sigma_1,\ldots,\sigma_m\}.
\]

\item The opposite-arrow conclusion for this \(\eta\) is the empty forward-fiber conclusion supplied in Step 1:
\[
R^{\mathrm{right}}_{m,K}(\Phi^{\mathrm{left}}_{m,K}(\eta))=\varnothing .
\]
The same root-polynomial argument also gives \(\sigma_i\ne0\) for every latent index \(i\).

\item For each \(i\), define \(\eta'\) by swapping the full reverse direct-source pair
\((\delta,(d_{m+1,r})_{r=2}^K)\) with the full \(i\)-th latent pair
\((\sigma_i,(d_{ir})_{r=2}^K)\), leaving all other coordinates fixed.  Since \(\eta\) is generic and \(\sigma_i\ne0\), this swapped point remains in
\(\Theta^{\mathrm{left},\circ}_{m,K}\).

\item This \(\eta'\) is, by construction, the reverse direct-latent swap required in the statement.

\item The reverse cumulant map is unchanged by the swap, again because the source sum is only reindexed by the corresponding transposition.  Thus
\[
\Phi^{\mathrm{left}}_{m,K}(\eta')
=
\Phi^{\mathrm{left}}_{m,K}(\eta).
\]

\item The swapped reverse parameter is not in the admissible \(G_m\)-orbit of \(\eta\): admissible relabelings keep the reverse direct coordinate fixed, while the constructed swap changes it from \(\delta\) to \(\sigma_i\), and genericity gives \(\delta\ne\sigma_i\).

\item Finally, if \(m\ge2\), the reverse same-arrow fiber dimension statement applied with the recovered multiset from Step 9 gives exact relative Zariski dimension
\[
\frac{m(m-1)}2
\]
for
\(R^{\mathrm{left}}_{m,K}(\Phi^{\mathrm{left}}_{m,K}(\eta))\).
\end{enumerate}
\end{proof}

\begin{proof}[Proof of \cref{obj:thm:exceptional-locus-codimension-one}]
Fix \(m\ge 1\) and set \(K=2m+2\).  Let
\[
D_m:=\frac{3(m+2)^2+(m+2)}{2}-5 .
\]
\begin{enumerate}
\item The geometric argument works in the affine space of the retained cumulants, those of orders \(2,\ldots,K\). In that space, the closure of the common-axis image is a proper irreducible component of \(\overline E_m\) and has dimension one less than each arrow image variety of \cref{obj:def:image-varieties}. The codimension claim follows from this dimension comparison and irreducibility. The preimage and low-dimensional claims then follow by unpacking \cref{obj:def:generic-full-fiber-compatibility}.

\item For the common-axis image, evaluate a \(D_m\times D_m\) Jacobian minor at an explicit common-axis parameter point. The matrix is block upper triangular: the ordinary-order block is Vandermonde-type, the top-order block is an augmented confluent block, and the lower-left block vanishes because cumulants of different orders have separate source-weight coordinates. Both diagonal determinants are nonzero. The Jacobian lower bound and the matching transcendence-degree upper bound therefore give dimension \(D_m\) for the common-axis image.

\item For the full forward arrow image, evaluate the analogous \((D_m+1)\times(D_m+1)\) Jacobian minor at an explicit forward parameter point. The same order separation makes the matrix block upper triangular, and both diagonal blocks have nonzero determinant. Hence the forward polynomial image has dimension at least \(D_m+1\); the coordinate-algebra upper bound gives equality.

\item Let \(C\) be the common-axis image closure in the retained cumulant coordinates, and let \(X^{\mathrm{right}}\) be the forward arrow image variety in the same coordinates. The inclusions \(C\subseteq \overline E_m\subseteq X^{\mathrm{right}}\) hold by construction. They are proper because a horizontal contraction polynomial vanishes on the reverse image but is nonzero at an explicit forward parameter point. Since \(X^{\mathrm{right}}\) is irreducible and the preceding dimensions are \(D_m\) and \(D_m+1\), no irreducible closed set can lie strictly between \(C\) and \(X^{\mathrm{right}}\). Exchanging the coordinates gives the same conclusion for the reverse image.

\item Projection to the retained cumulant coordinates is injective on the truncated objects under consideration, so the dimension comparison transfers to the ambient cumulant space. The forward and reverse arrow image varieties are irreducible polynomial-image closures, and \(\overline E_m\) is contained in both. On the common-axis component, the horizontal or vertical contraction polynomial defines a principal component and rules out an intermediate irreducible set between that component and the corresponding arrow variety. Thus \(\overline E_m\) has codimension exactly \(1\) in both \(C^{\mathrm{right}}_{m,K}\) and \(C^{\mathrm{left}}_{m,K}\).

\item If \(m\ge 2\), exact codimension \(1\) excludes exact codimension \(m\): a longer irreducible chain would contain an intermediate closed set between the common-axis component and the ambient arrow variety, contradicting the preceding conclusion.

\item For the forward preimage, \cref{obj:def:generic-full-fiber-compatibility} identifies \(H^{\mathrm{right}}_m\) with the parameters whose forward image lies in the generic full-fiber compatibility locus.  Unpacking this definition, and using \cref{obj:def:generic-parameter-loci} and \cref{obj:def:fiber-correspondences}, gives exactly
\[
\theta\in H^{\mathrm{right}}_m
\quad\Longleftrightarrow\quad
\theta\in\Theta^{\mathrm{right},\circ}_{m,K}
\ \text{and}\
R^{\mathrm{left}}_{m,K}\!\left(\Phi^{\mathrm{right}}_{m,K}(\theta)\right)\neq\varnothing .
\]
Coordinates outside orders \(2,\ldots,K\) vanish under the truncation convention.

\item The reverse preimage is identical with the two arrows interchanged.  Thus \cref{obj:def:generic-full-fiber-compatibility}, \cref{obj:def:generic-parameter-loci}, and \cref{obj:def:fiber-correspondences} give
\[
\eta\in H^{\mathrm{left}}_m
\quad\Longleftrightarrow\quad
\eta\in\Theta^{\mathrm{left},\circ}_{m,K}
\ \text{and}\
R^{\mathrm{right}}_{m,K}\!\left(\Phi^{\mathrm{left}}_{m,K}(\eta)\right)\neq\varnothing .
\]

\item Finally, for \(m=1\) and \(m=2\), the worked systems of \cref{obj:def:worked-compatibility-instances} are exactly the retained-coordinate incidence descriptions of the generic compatibility locus. A point of \(E_m\) supplies forward and reverse parameters with equal retained cumulants and hence a solution of \(A_m\). Conversely, a solution \(p=(p^{\mathrm{right}},p^{\mathrm{left}})\) supplies parameters in the two fibers with the required generic side. Coordinates beyond the truncation order vanish by convention. Taking \(m=1\) and \(m=2\) gives the displayed formulas with \(A_1\) and \(A_2\).
\end{enumerate}
\end{proof}

\begin{proof}[Proof of \cref{obj:thm:improved-real-information-order}]
Fix \(m\ge 3\), and put \(L=2m+1\).

\begin{enumerate}
\item The separation input at this order is obtained by the lower-order apolar
calculation.  On the forward parameter space, let
\[
\Delta^{\mathrm{right}}_m
  =
  \Bigl(\prod_{i=1}^{m}\rho_i\Bigr)\cdot
  \bigl(\text{the forward lower-order contraction-rank determinant}\bigr),
\]
and on the reverse parameter space let
\[
\Delta^{\mathrm{left}}_m
  =
  \Bigl(\prod_{i=1}^{m}\sigma_i\Bigr)\cdot
  \bigl(\text{the reverse lower-order contraction-rank determinant}\bigr).
\]
Evaluation at explicit parameter points shows that both determinant factors are
nonzero polynomials; hence \(\Delta^{\mathrm{right}}_m\) and
\(\Delta^{\mathrm{left}}_m\) are nonzero real polynomials.  Define
\[
E^{\mathrm{right}}
  :=\{\theta:\Delta^{\mathrm{right}}_m(\theta)=0\},
\qquad
E^{\mathrm{left}}
  :=\{\eta:\Delta^{\mathrm{left}}_m(\eta)=0\}.
\]
Thus both sets are proper real algebraic subsets of their respective real
parameter spaces.  The truncated-moment interior fact at order \(L\), together
with these explicit nonvanishing evaluations, gives feasible real points outside each
zero locus:
\[
F^{\mathrm{right}}_{m,L}\setminus E^{\mathrm{right}}\ne\varnothing,
\qquad
F^{\mathrm{left}}_{m,L}\setminus E^{\mathrm{left}}\ne\varnothing .
\]

\item Let \(\theta\in F^{\mathrm{right}}_{m,L}\setminus E^{\mathrm{right}}\).
After viewing \(\theta\) as a complex parameter, the nonvanishing of
\(\Delta^{\mathrm{right}}_m(\theta)\) gives all \(\rho_i\ne0\) and
nonvanishing of the forward contraction-rank determinant.  Feasibility gives
\(\gamma\ne0\) and the required distinctness of the forward loading slopes, so
the forward loading directions used in the lower-order apolar blocks are
injectively indexed and, together with the nonzero \(\rho_i\)'s, have the needed
nonzero second coordinates.  The determinant condition gives injectivity of the
associated lower-order weighted contraction map.  Suppose, toward a
contradiction, that some \(\eta\in F^{\mathrm{left}}_{m,L}\) satisfies
\[
\Phi^{\mathrm{right}}_{m,L}(\theta)
 =
\Phi^{\mathrm{left}}_{m,L}(\eta).
\]
The equality remains true after complexification, since the cumulant maps of
\cref{obj:def:forward-cumulant-map} and
\cref{obj:def:reverse-cumulant-map} have real polynomial coordinates.

\item Let \(Q_{\mathrm{left}}\) be the degree-\(m+2\) support-annihilator
polynomial of the reverse loading directions, and let \(Q_{\mathrm{right}}\) be
the corresponding support-annihilator polynomial of the forward loading
directions, in the notation of \cref{obj:def:apolar-notation}.
The polynomial \(Q_{\mathrm{left}}\) annihilates the reverse divided-power
cumulant blocks in orders \(m+2,\ldots,2m+1\); by the assumed equality of
cumulants, it also annihilates the corresponding forward blocks.  The forward
lower-order apolar kernel identity, applied with the slope and rank conditions
from the previous step, therefore yields a scalar \(c\) such that
\[
Q_{\mathrm{left}}=c\,Q_{\mathrm{right}}.
\]
But every reverse support annihilator is divisible by the horizontal-axis
factor, whereas the forward support annihilator is not divisible by that factor
when \(\gamma\ne0\) and all \(\rho_i\ne0\).  If \(c=0\), then
\(Q_{\mathrm{left}}=0\), contradicting nonzero support annihilation.  If
\(c\ne0\), divisibility transfers from \(Q_{\mathrm{left}}\) to
\(Q_{\mathrm{right}}\), again a contradiction.  Hence no such reverse
\(\eta\) exists.

\item The reverse argument is identical with the roles interchanged.  For
\(\eta\in F^{\mathrm{left}}_{m,L}\setminus E^{\mathrm{left}}\), the
nonvanishing of \(\Delta^{\mathrm{left}}_m(\eta)\) gives all
\(\sigma_i\ne0\) and nonvanishing of the reverse contraction-rank determinant.
Feasibility gives \(\delta\ne0\) and injectivity of the relevant reverse loading
slopes; together with the nonzero \(\sigma_i\)'s this gives the needed nonzero
first coordinates.  The determinant condition gives injectivity of the reverse
lower-order weighted contraction map.  If
\[
\Phi^{\mathrm{left}}_{m,L}(\eta)
 =
\Phi^{\mathrm{right}}_{m,L}(\theta)
\]
for some \(\theta\in F^{\mathrm{right}}_{m,L}\), then the forward support
annihilator annihilates the reverse lower-order blocks and must be a scalar
multiple of the reverse support annihilator.  The forward support annihilator is
divisible by the vertical-axis factor, while the reverse one is not divisible by
that factor under \(\delta\ne0\) and all \(\sigma_i\ne0\).  The same scalar
case split gives a contradiction.  Thus no such forward \(\theta\) exists.

\item The preceding four steps prove generic real separation at order
\(L=2m+1\), with the theorem-local exceptional sets
\(E^{\mathrm{right}}\) and \(E^{\mathrm{left}}\).  Since \(m\ge3\), also
\(2\le 2m+1\).  Therefore \(2m+1\) is an admissible element in the set whose
infimum defines \(K^\star(m)\) in \cref{obj:def:information-order},
and hence
\[
K^\star(m)\le 2m+1 .
\]

\item Finally, if \(K^\star(m)=2m+2\), the just-proved inequality would give
\(2m+2\le 2m+1\) after comparing the corresponding finite elements of
\(\mathbb N_\infty\), which is impossible.  Hence
\[
K^\star(m)\ne 2m+2 .
\]

\end{enumerate}
\end{proof}

\begin{proof}[Proof of \cref{obj:thm:generic-apolar-arrow-recovery}]
Put \(K=2m+2\) and \(n=m+2\).  We argue in nine steps.

\begin{enumerate}
\item The truncated Hamburger moment interior theorem gives, at order \(K\), a nonempty Euclidean open set of admissible centered non-Gaussian one-dimensional source cumulant vectors, realized by laws with finite \(K\)-th moment.  This is the real feasibility input used below to make prescribed nonvanishing polynomial conditions hold inside the feasible region of \cref{obj:def:real-feasible-regions}.%

\item On the forward side, choose a nonzero determinant polynomial \(M_{\mathrm r}\) whose nonvanishing makes the forward weighted contraction map injective; on the reverse side choose the analogous nonzero determinant polynomial \(M_{\mathrm l}\).  In each case the determinant is nonzero at an explicit parameter whose source-cumulant coordinates vanish outside the orders \(2,\ldots,K\), so the same determinant remains nonzero after restricting to the truncated coordinate chart.%

\item Let \(\Delta_{\mathrm{gen}}\) be the polynomial cutting out the complement of the generic locus of \cref{obj:def:generic-parameter-loci}: it is the product of the direct slope, all direct-versus-latent slope differences, all latent slope differences, and all retained source-cumulant coordinates.  Let
\[
\Delta_\rho=\prod_{i=1}^m \rho_i
\]
in the forward chart, and use the identical product in the reverse chart with the reverse finite slopes.  Define
\[
R_{\mathrm r}:=\Delta_{\mathrm{gen}}\Delta_\rho M_{\mathrm r},
\qquad
R_{\mathrm l}:=\Delta_{\mathrm{gen}}\Delta_\rho M_{\mathrm l}.
\]
The preceding step and the explicit nonzero factors show that both \(R_{\mathrm r}\) and \(R_{\mathrm l}\) are nonzero on the retained-coordinate parameter spaces.%

\item Define \(U_m^{\mathrm{right}}\) to be the truncated forward parameter space cut by \(R_{\mathrm r}\ne0\), and \(U_m^{\mathrm{left}}\) to be the truncated reverse parameter space cut by \(R_{\mathrm l}\ne0\).  A principal nonvanishing locus whose truncated restriction is not the zero polynomial is relatively Zariski open and relatively Zariski dense.  Since \(R_{\mathrm r}\) and \(R_{\mathrm l}\) contain \(\Delta_{\mathrm{gen}}\), these loci are contained in the corresponding generic parameter loci; since they contain \(\Delta_\rho\), all finite latent slopes are nonzero on them; and since they contain the determinant factors, the corresponding weighted contraction maps are injective.%

\item Set
\[
O^{\mathrm{right}}=\{\theta\in\Theta^{\mathrm{right}}_{m,K}(\mathbb R):
R_{\mathrm r}(\theta_{\mathbb C})\ne0\},
\qquad
O^{\mathrm{left}}=\{\eta\in\Theta^{\mathrm{left}}_{m,K}(\mathbb R):
R_{\mathrm l}(\eta_{\mathbb C})\ne0\}.
\]
These are Euclidean open by continuity.  Applying the moment-interior nonvanishing statement to \(R_{\mathrm r}\) and \(R_{\mathrm l}\) gives feasible real points in each open set.  If a feasible real parameter lies in the appropriate open set, then its coordinatewise complexification satisfies the truncated support equations and the same nonvanishing condition, hence belongs to the corresponding \(U\)-locus.%

\item Fix \(\theta\in U_m^{\mathrm{right}}\).  From the construction, \(\theta\) is generic, \(\gamma\ne0\), all \(\rho_i\ne0\), the finite forward loading slopes are pairwise distinct and nonzero, and the forward weighted contraction map is injective.  Suppose that a reverse parameter \(\eta\) had the same truncated cumulants.  Its reverse support-annihilator is homogeneous of degree \(n\), annihilates the reverse divided-power blocks, and therefore, by equality of cumulants, annihilates the forward blocks.  The forward apolar kernel identity then forces it to be a scalar multiple of the forward support-annihilator.  The scalar is nonzero because a support-annihilator of the loading directions is nonzero.  But the reverse support-annihilator is divisible by \(y\), whereas the forward support-annihilator is not divisible by \(y\), since \(\gamma\) and all \(\rho_i\) are nonzero.  This contradiction proves
\[
R^{\mathrm{left}}_{m,K}\!\left(\Phi^{\mathrm{right}}_{m,K}(\theta)\right)=\varnothing .
\]%

\item For the same \(\theta\), take
\[
Q(z)=\prod_{r\in\{\gamma,\rho_1,\ldots,\rho_m\}}(z-r),
\]
with the multiset interpretation of the product.  Genericity makes the entries \(\gamma,\rho_1,\ldots,\rho_m\) pairwise distinct, and the extra factor \(\Delta_\rho\) together with genericity excludes \(0\); hence \(Q\ne0\), \(Q\) is squarefree, \(Q(0)\ne0\), and its root multiset is exactly \(\{\gamma,\rho_1,\ldots,\rho_m\}\).  If another generic forward parameter \(\theta''\) has the same forward cumulants, its own support-annihilator lies in the same forward apolar kernel; the kernel identity makes it a nonzero scalar multiple of the support-annihilator of \(\theta\), and dehomogenizing in the \(x\)-chart gives the same root multiset.%

\item Let \(D\) be the forward loading-direction support and let \(Q_D\) be its homogeneous support-annihilator as in \cref{obj:def:apolar-notation}.  Each loading direction has at least one nonzero coordinate, so \(Q_D\ne0\).  The fixed vertical loading direction contributes the factor \(x\), while the nonzero finite slopes and \(\gamma\ne0\) exclude a factor \(y\).  Finally, the forward apolar kernel identity gives, for every homogeneous binary form \(q\) of degree \(n\),
\[
\bigl[q(\partial)f_{n+k}=0\ \text{for all }0\le k\le m\bigr]
\Longleftrightarrow
\exists c\in\mathbb C:\ q=cQ_D .
\]%

\item The reverse proof is the same argument with \(x\) and \(y\), and right and left, interchanged.  For \(\eta\in U_m^{\mathrm{left}}\), genericity and \(R_{\mathrm l}\ne0\) give \(\delta\ne0\), all \(\sigma_i\ne0\), distinct nonzero reverse finite loading coordinates, and injectivity of the reverse weighted contraction map.  A hypothetical forward parameter with the same cumulants would have a forward support-annihilator lying in the reverse apolar kernel; hence it would be a nonzero scalar multiple of the reverse support-annihilator.  This contradicts divisibility by \(x\), because the forward support-annihilator is divisible by \(x\) and the reverse one is not.  Thus the forward fiber is empty.  The univariate polynomial
\[
Q(z)=\prod_{r\in\{\delta,\sigma_1,\ldots,\sigma_m\}}(z-r)
\]
is nonzero, squarefree, nonzero at \(0\), and has the stated root multiset; the same dehomogenized-kernel argument shows that every generic reverse parameter with the same cumulants has that multiset.  For the reverse loading-direction support \(D\), the homogeneous \(Q_D\) is nonzero, is divisible by \(y\), is not divisible by \(x\), and its degree-\(n\) common apolar kernel is exactly the line \(\mathbb C Q_D\).%

\end{enumerate}
\end{proof}

\bibliographystyle{plainnat}

\bibliography{references}

\end{document}
